\chapter{Tutorials: the Solid Engine (msg-solid)}
\label{ch:solid}

The \texttt{msg-solid} engine reduces a \emph{solid} structure gene (SG) --- a
mesh of ordinary continuum elements in one, two or three dimensions --- to a
macroscopic constitutive law, and then runs the same construction backwards to
recover the pointwise three-dimensional strain and stress inside that gene.
Everything in this chapter is reached through two entry points.

\begin{table}[htbp]
\centering
\footnotesize
\begin{tabular}{llp{4.8cm}}
\toprule
entry point & module & structure gene it reads \\
\midrule
\texttt{rm\_plate\_msg} & \texttt{opensg\_solid.rm\_plate\_1D.msg\_rm\_plate}
  & a 1-D line mesh through a wall thickness \\
\texttt{plate\_homo\_2d} & \texttt{opensg\_solid.sg\_homo}
  & a 1-D, 2-D or 3-D solid SG \\
\bottomrule
\end{tabular}
\caption{The two homogenization entry points of \texttt{msg-solid}.}
\label{tab:solid-entry}
\end{table}

The first is the dedicated through-thickness (laminate) route: the wall is a
stack of plies, nothing varies in-plane, so the gene is one-dimensional and the
result is a plate law. The second is the general route. Its \texttt{n\_model}
argument selects which macro model the same SG is reduced to.

\begin{table}[htbp]
\centering
\footnotesize
\begin{tabular}{clll}
\toprule
\texttt{n\_model} & macro model & row and column order & title infix \\
\midrule
1 & beam & \texttt{[eps11 gam12 gam13 k1 k2 k3]} & \texttt{Timoshenko Beam} \\
2 & plate & \texttt{[N11 N22 N12 M11 M22 M12]} & \texttt{Classical Plate} \\
3 & 3-D elastic & Voigt \texttt{[11 22 33 23 13 12]} & \texttt{Cauchy Continuum} \\
\bottomrule
\end{tabular}
\caption{\texttt{n\_model} on \texttt{plate\_homo\_2d}. The infix goes into
\texttt{The Effective <infix> Stiffness Matrix} and is always the macro law's
console title, so the file and the terminal read the same. Two of the five
titles belong to the refinement switch rather than to \texttt{n\_model}: the
classical beam $4\times4$ is \texttt{Euler-Bernoulli Beam}, and the plate route
becomes \texttt{Reissner-Mindlin} when the shear-refined $8\times8$ is
produced.}
\label{tab:solid-nmodel}
\end{table}

Every \texttt{plate\_homo\_2d} call writes a timed SwiftComp-\texttt{.K}-layout
report --- effective stiffness, effective compliance, and, for the 3-D solid law
only, an engineering-constants block --- opening with an \texttt{ OpenSG} banner
and closing with a \texttt{ Time taken: } line in seconds. Unless you pass
\texttt{plot=False} it also writes a \texttt{<base>\_mesh.png} with the elements
coloured by material. One rule governs the names of both: inside
\texttt{plate\_homo\_2d} the output base is
\texttt{base = os.path.splitext(sc\_path)[0]}, so the report is
\emph{always} \texttt{<input basename>.out}, regardless of what the script calls
the case.

The eight runnable examples live under \texttt{examples/OpenSG-solid}. Each
runner is meant to be executed from inside its own folder, because the SG file
is named relative to the working directory.

\begin{table}[htbp]
\centering
\footnotesize
\begin{tabular}{lp{3.2cm}p{5.4cm}}
\toprule
folder & script & what it does \\
\midrule
\texttt{1\_get\_plate\_props\_from\_1DSG} & \texttt{1d\_sg.py} & layup $\rightarrow$ 1-D SG $\rightarrow$ plate law \\
\texttt{2\_get\_plate\_stress\_from\_1DSG} & \texttt{rm\_dehom.py}, \texttt{rm\_dehom\_ff.py} & plate resultants $\rightarrow$ 3-D ply stress \\
\texttt{3\_get\_plate\_props\_from\_2D\_SG} & \texttt{plate\_homo\_2dsg.py} & 2-D cell SG $\rightarrow$ plate ABD \\
\texttt{4\_get\_plate\_dehom\_from\_2DSG} & \texttt{plate\_dehom\_2dsg.py} & macro plate strain $\rightarrow$ fields in the cell \\
\texttt{5\_get\_solid\_props\_from\_SG} & \texttt{solid\_homo\_2dsg.py} & 2-D cell SG $\rightarrow$ equivalent 3-D solid \\
\texttt{6\_get\_solid\_props\_from\_3D\_SG} & \texttt{tpms\_solid\_props.py} and others & 3-D SG $\rightarrow$ equivalent 3-D solid \\
\texttt{7\_get\_beam\_props\_from\_SG} & \texttt{beam\_homo\_sg.py} & 2-D cross-section SG $\rightarrow$ Timoshenko $6\times6$ \\
\texttt{8\_get\_beam\_dehom\_from\_SG} & \texttt{beam\_dehom\_sg.py} & macro beam strain $\rightarrow$ fields in the section \\
\bottomrule
\end{tabular}
\caption{The eight \texttt{msg-solid} example folders.}
\label{tab:solid-examples}
\end{table}

A note on what you will find in those folders before you run anything. Some of
the files named as outputs below are regenerated on every run and are simply not
checked in; where that is the case it is said explicitly. Nothing in this
chapter depends on a pre-computed file being present.

\section{Plate properties from a through-thickness 1-D SG}
\label{sec:solid-plate1d}

\subsection{What it computes}

The structure gene is the wall discretised through its own thickness: one
element per ply, quartic by default, with each node carrying the three warping
(fluctuation) degrees of freedom. Minimising the 3-D strain energy over that
line mesh order by order in $h/\ell$ gives the warping ladder --- the
zeroth-order columns $V_0$ and the first-gradient columns $V_{11}$, $V_{12}$ ---
from which the classical plate stiffness follows as
%
\begin{equation}
\mathbf{A}_6 = \big\langle (\mathbf{\Gamma}_\varepsilon
 + \mathbf{\Gamma}_h \mathbf{V}_0)^{\mathsf T}\,
 \mathbf{C}(z)\,
 (\mathbf{\Gamma}_\varepsilon + \mathbf{\Gamma}_h \mathbf{V}_0) \big\rangle ,
\label{eq:solid-A6}
\end{equation}
%
where $\langle\cdot\rangle$ is the through-thickness integral and
$\mathbf{C}(z)$ the ply stiffness rotated by its fibre angle. With
\texttt{model: 1} the run additionally identifies the $2\times2$ MSG
transverse-shear block $\mathbf{G}$ from the second-order energy and assembles
the $8\times8$ law
%
\begin{equation}
\mathbf{ABDG} =
\begin{bmatrix} \mathbf{A}_6 & \mathbf{0} \\ \mathbf{0} & \mathbf{G} \end{bmatrix},
\label{eq:solid-abdg}
\end{equation}
%
in which the transverse shears are uncoupled from the membrane and bending
measures. Row and column order is
\texttt{[e11 e22 g12 k11 k22 k12 2g13 2g23]}, and that string is printed in the
banner of every \texttt{.out} so a file can always be read without the manual.

The whole capability is one file you edit and five functions that run. The SG
mesh is written to disk and read back before the solve, so the mesh file --- not
the database --- is what the homogenization actually sees.

\captree{
  \node[inputnode] (y) {\texttt{layup\_db.yaml}\\
    \texttt{model}, \texttt{fraction}, \texttt{mesh},\\
    \texttt{materials}, \texttt{layup}};
  \node[funcnode, below=of y] (f1) {\texttt{segment\_plate.plate\_sg\_yaml}\\
    the through-thickness node grid,\\
    \texttt{n\_per\_layer} elements per ply};
  \node[outnode, right=of f1] (o1) {\texttt{1dsg.yaml}\\ \texttt{1dsg.png}};
  \node[funcnode, below=of f1] (f2)
    {\texttt{segment\_plate.read\_plate\_sg\_yaml}\\
     the mesh back as thick, angles,\\ mat\_names, material\_db};
  \node[funcnode, below=of f2] (f3) {\texttt{msg\_rm\_plate.rm\_plate\_msg}\\
    warping ladder \texttt{V0}, \texttt{V11}, \texttt{V12}};
  \node[datanode, below=of f3] (d1) {\texttt{r["A6"]} $6\times6$ ABD\\
    \texttt{r["ABDG"]} $8\times8$ at \texttt{model: 1}};
  \node[funcnode, below=of d1] (f4) {\texttt{sg\_homo.write\_sc\_K}\\
    stiffness, compliance, banner, timing};
  \node[outnode, right=of f4] (o2) {\texttt{layup\_db\_plate\_homo.out}};
  \draw[capedge] (y) -- (f1);
  \draw[capedge] (f1) -- (o1);
  \draw[capedge] (o1) -- (f2);
  \draw[capedge] (f2) -- (f3);
  \draw[capedge] (f3) -- (d1);
  \draw[capedge] (d1) -- (f4);
  \draw[capedge] (f4) -- (o2);
}

\subsection{The input you edit}

Everything lives in \texttt{layup\_db.yaml}. The script reads that file and
nothing else, so a new laminate means editing the YAML, never the code.

\begin{table}[htbp]
\centering
\footnotesize
\begin{tabular}{lp{7.2cm}p{2.6cm}}
\toprule
key & what it sets & value in the shipped file \\
\midrule
\texttt{model} & \texttt{0} = classical $6\times6$ ABD, \texttt{1} = $8\times8$ ABDG & \texttt{1} \\
\texttt{fraction} & reference surface as a fraction of thickness & \texttt{0.5} \\
\texttt{mesh: n\_per\_layer} & SG elements per physical ply & \texttt{1} \\
\texttt{mesh: elem\_order} & element order; \texttt{4} = 5-noded quartic & \texttt{4} \\
\texttt{materials} & per material \texttt{E: [E1,E2,E3]}, \texttt{G: [G12,G13,G23]}, \texttt{nu: [nu12,nu13,nu23]}, \texttt{rho} & \texttt{ge}, \texttt{herex} \\
\texttt{layup} & one entry per ply, bottom ply first & 9 plies \\
\bottomrule
\end{tabular}
\caption{The six keys \texttt{1d\_sg.py} reads from \texttt{layup\_db.yaml}.}
\label{tab:solid-layupdb}
\end{table}

For \texttt{fraction}, \texttt{0} is the bottom (OML) face, \texttt{0.5} the
mid-surface and \texttt{1} the top face. It is a property of the \emph{mesh},
not of the homogenization call: the generated node coordinates are measured from
the reference plane, so they run from $-\text{fraction}\cdot h$ to
$(1-\text{fraction})\cdot h$ and $x_3 = 0$ \emph{is} the reference. The $B$
block of the ABD depends on this choice, so homogenize, analyze and recover with
the same \texttt{fraction}.

The shipped case is a $[0/90/0/90/\text{core}]_s$ sandwich: eight
graphite/epoxy plies of $0.9525$~mm ($E_1 = 128$~GPa, $E_2 = E_3 = 11$~GPa,
$G_{12}=G_{13}=4.48$~GPa, $G_{23}=1.53$~GPa, $\rho = 1500$~kg/m$^3$) around a
$144.78$~mm HEREX C70.130 PVC foam core ($E = 103.63$~MPa, $G = 50$~MPa,
$\nu = 0.32$, $\rho = 130$~kg/m$^3$), total $h = 0.1524$~m.

The same \texttt{layup\_db.yaml} also carries a \texttt{V2} flag, a
\texttt{plate:} block and a \texttt{dehom:} block. Those are read by the
recovery scripts of Section~\ref{sec:solid-dehom1d}; \texttt{1d\_sg.py} itself
uses only the six keys in Table~\ref{tab:solid-layupdb}.

\subsection{Run it}

\begin{lstlisting}[style=shellst]
cd examples/OpenSG-solid/1_get_plate_props_from_1DSG
python 1d_sg.py
\end{lstlisting}

A different database file may be passed as the single positional argument:

\begin{lstlisting}[style=shellst]
python 1d_sg.py my_other_layup.yaml
\end{lstlisting}

\subsection{What comes out}

\begin{table}[htbp]
\centering
\small
\begin{tabular}{lp{9.6cm}}
\toprule
file & content \\
\midrule
\texttt{1dsg.yaml} & the generated 1-D SG mesh: an \texttt{sg:} header
  (\texttt{type: plate\_1d}, \texttt{elem\_order: 4}, \texttt{n\_per\_layer: 1},
  \texttt{reference\_fraction: 0.5}, \texttt{thickness}, \texttt{n\_ply}), 37
  nodes given as single $x_3$ coordinates from $-0.0762$ to $+0.0762$~m, 9
  five-noded elements, the material blocks with densities, and one element set
  plus one \texttt{sections:} entry per ply carrying its material, angle and
  thickness \\
\texttt{1dsg.png} & the mesh figure, plies coloured by material \\
\texttt{layup\_db\_plate\_homo.out} & the plate law, named
  \texttt{<layup\_db stem>\_plate\_homo.out} \\
\bottomrule
\end{tabular}
\caption{Outputs of \texttt{1d\_sg.py}.}
\label{tab:solid-ex1out}
\end{table}

The \texttt{.out} opens with a banner that records the whole run, then the
matrix whose title follows \texttt{model} --- \texttt{Reissner-Mindlin}
here, \texttt{Classical Plate} at \texttt{model: 0}:

\begin{lstlisting}[style=shellst,basicstyle=\ttfamily\tiny,breaklines=true]
 OpenSG msg-solid plate model from 1dsg.yaml: 9 plies, h = 0.152400 m, fraction = 0.5, rho*h = 30.251400 kg/m^2, rows [e11 e22 g12 k11 k22 k12 2g13 2g23]

 The Effective Reissner-Mindlin Stiffness Matrix
 --------------------------------------------
     5.4916502E+008     2.6417019E+007     1.5524611E-010     3.6182957E-006     3.6122933E-006     1.0244011E-024     0.0000000E+000     0.0000000E+000
\end{lstlisting}

The effective compliance follows in the same layout and the file closes with
\texttt{ Time taken: 0.82 sec}. Engineering constants are printed only for the
3-D solid law, never for a plate law, because a plate stiffness is not a
material stiffness.

\subsection{What the numbers mean}

The eight diagonal entries of the shipped run, in SI (the code is unit-agnostic
and simply follows the input, which is metres and pascals here):

\begin{table}[htbp]
\centering
\small
\begin{tabular}{lll}
\toprule
term & value & units \\
\midrule
$A_{11}$ & $5.4916502\times10^{8}$ & N/m \\
$A_{22}$ & $5.4916502\times10^{8}$ & N/m \\
$A_{66}$ & $4.1376600\times10^{7}$ & N/m \\
$D_{11}$ & $3.0005458\times10^{6}$ & N$\cdot$m \\
$D_{22}$ & $2.9371144\times10^{6}$ & N$\cdot$m \\
$D_{66}$ & $2.0111708\times10^{5}$ & N$\cdot$m \\
$G_{11}$ & $7.7009772\times10^{6}$ & N/m \\
$G_{22}$ & $7.5423712\times10^{6}$ & N/m \\
\bottomrule
\end{tabular}
\caption{Diagonal of the $8\times8$ plate law of the shipped sandwich.}
\label{tab:solid-ex1diag}
\end{table}

Three checks are worth making on any run of your own, and all three are visible
in this one.

\begin{itemize}
\item \textbf{The section mass.} The banner's $\rho h = 30.2514$~kg/m$^2$ equals
  $8 \times 0.0009525 \times 1500 + 0.14478 \times 130
  = 11.43 + 18.8214$~kg/m$^2$, so the layup that reached the solver is the
  layup you wrote.
\item \textbf{The $B$ block.} The stack is symmetric about the mid-surface and
  \texttt{fraction: 0.5} puts the reference plane there, so membrane-bending
  coupling must vanish. It does: the largest entry of the $B$ block is
  $3.6183\times10^{-6}$ against an $A_{11}$ of $5.4917\times10^{8}$, fourteen
  orders down.
\item \textbf{The transverse-shear block.} $G_{11} = 7.7010\times10^{6}$~N/m
  sits $6.4\,\%$ above the crude core-only estimate
  $G_{\rm core} h_{\rm core} = 50\times10^{6} \times 0.14478
  = 7.239\times10^{6}$~N/m, which is the expected magnitude for a sandwich whose
  shear is carried almost entirely by the foam. $A_{11}$, by contrast, is set by
  the stiff faces.
\end{itemize}

\begin{figure}[htbp]
\centering
\includegraphics[width=0.5\textwidth]{\FIGDIR/plate_1dsg.png}
\caption{The through-thickness 1-D structure gene of the sandwich laminate, as
written by the run to \texttt{1dsg.png}. The gene is a line, drawn with
thickness for legibility: the thin graphite/epoxy face plies are at top and
bottom, the foam core between them, and the dashed line marks the reference
surface at mid-thickness.}
\label{fig:solid-plate1dsg}
\end{figure}

\subsection{What it is validated against}

\texttt{src/opensg\_solid/rm\_plate\_1D/tests/test\_msg\_rm\_plate.py} gates the
construction on analytic anchors and on an independent plain-NumPy re-assembly
of the same SG matrices. For an isotropic plate with $\nu = 0$ the
transverse-shear block comes out exactly $\tfrac56 G h$, the textbook
shear-correction factor, with the least-squares residual $U^*$ at machine zero
(\texttt{test\_isotropic\_nu0\_gives\_five\_sixths}); $\mathbf{A}_6$ equals the
classical lamination-theory \texttt{msg\_materials.compute\_ABD\_matrix} at the
matching reference plane (\texttt{test\_A6\_matches\_compute\_ABD});
$\mathbf{A}_6(\text{fraction})$ equals the offset transform of
$\mathbf{A}_6(0)$ (\texttt{test\_reference\_plane\_transform}); and the warping
fields satisfy the Euler-Lagrange equations and the gauge condition
$\psi^{\mathsf T} H V = 0$ for $V_0$, $V_{11}$ and $V_{12}$ directly
(\texttt{test\_euler\_lagrange\_residuals}, \texttt{test\_H\_identities\_and\_gauge}).
The formulation is that of Yu, Hodges and Volovoi, \emph{Computers \& Structures}
81 (2003) 439--454.

\section{Through-thickness stress recovery}
\label{sec:solid-dehom1d}

\subsection{What it computes}

Homogenization threw information away; recovery puts it back. Given the section
resultants at a point of the plate, the recovery evaluates the 3-D strain,
stress, warping and displacement at every depth of the gene. The construction is
linear in \textbf{42 drivers},
%
\begin{equation}
\mathbf{D} = [\,\mathbf{E}_6,\ \mathbf{E}_{6,1},\ \mathbf{E}_{6,2},\
 \mathbf{E}_{6,11},\ \mathbf{E}_{6,12},\ \mathbf{E}_{6,22},\ \mathbf{qt}_6\,],
\label{eq:solid-42}
\end{equation}
%
so the code builds three operators once per depth from 42 unit evaluations ---
$\mathbf{T}_E,\mathbf{T}_S \in \mathbb{R}^{6\times42}$ and
$\mathbf{T}_W \in \mathbb{R}^{3\times42}$ --- after which every station is one
matrix-vector product, \texttt{jax.vmap}-ed over the whole field. The plate
strains come from the section law itself,
$\mathbf{E}_6 = \mathbf{A}_6^{-1}\mathbf{F}_6$, so a recovery never bypasses the
homogenization.

Both drivers of this section --- the single station of \texttt{rm\_dehom.py} and
the whole field of \texttt{rm\_dehom\_ff.py} --- walk the same core functions in
\texttt{opensg\_solid.rm\_plate\_1D.rm\_dehom}. They differ only in where the
42-driver vector comes from.

\captree{
  \node[inputnode] (i1) {\texttt{1dsg.yaml}\\
    the 1-D SG of Section~\ref{sec:solid-plate1d}};
  \node[funcnode, below=of i1] (f1) {\texttt{rm\_dehom.load\_sg}\\
    \texttt{read\_plate\_sg\_yaml}, then\\
    \texttt{rm\_plate\_msg} re-run;\\
    \texttt{S6} = inverse of \texttt{A6}};
  \node[inputnode, right=of f1] (i2) {\texttt{F6}, \texttt{Q2v} in the USER\\
    INPUT block, \emph{or} a\\ \texttt{.ff} station table plus the\\
    \texttt{dehom:} block of \texttt{layup\_db.yaml}};
  \node[funcnode, below=of f1] (f2) {\texttt{rm\_dehom.z\_stations}\\
    \texttt{rm\_dehom.build\_ops}\\
    \texttt{TE}, \texttt{TS} $(n_z,6,42)$, \texttt{TW} $(n_z,3,42)$};
  \node[funcnode, below=of f2] (f3) {\texttt{rm\_dehom.assemble\_driver},\\
    or \texttt{np.gradient} on a lattice /\\
    \texttt{rm\_dehom.mls\_gradients} on a cloud};
  \node[datanode, below=of f3] (d1) {\texttt{Eps}, \texttt{Sig}, \texttt{Warp}
    per depth\\ then the \texttt{Q1}/\texttt{Q2} rescale};
  \node[funcnode, below=of d1] (f4)
    {\texttt{rm\_dehom.material\_rotations}\\ laminate frame to ply frame};
  \node[funcnode, below=of f4] (f5) {\texttt{rm\_dehom.write\_field}\\
    \texttt{rm\_dehom.write\_vtk}};
  \node[outnode, right=of f5] (o1) {\texttt{1dsg\_dehom.SM/.EM/.U/.vtk}\\
    or \texttt{<ff stem>\_ff.SM/.EM/.U}\\ and \texttt{<ff stem>\_ff\_gauss.vtk}};
  \draw[capedge] (i1) -- (f1);
  \draw[capedge] (f1) -- (f2);
  \draw[capedge] (f2) -- (f3);
  \draw[capedge] (i2) -- (f3);
  \draw[capedge] (f3) -- (d1);
  \draw[capedge] (d1) -- (f4);
  \draw[capedge] (f4) -- (f5);
  \draw[capedge] (f5) -- (o1);
}

Which driver feeds which component is the thing to understand before choosing an
input.

\begin{table}[htbp]
\centering
\footnotesize
\begin{tabular}{p{2.6cm}p{3.6cm}p{4.2cm}p{3.0cm}}
\toprule
driver & where it comes from & what it produces & without it \\
\midrule
$\mathbf{E}_6$ & $\mathbf{A}_6^{-1}\mathbf{F}_6$ at the station
  & $\sigma_{11}$, $\sigma_{22}$, $\sigma_{12}$ layer by layer & nothing \\
$\mathbf{E}_{6,1}$, $\mathbf{E}_{6,2}$ & in-plane gradients of the resultants
  & $\sigma_{13}$, $\sigma_{23}$ (the $V_{11}/V_{12}$ warping)
  & $\sigma_{13}=\sigma_{23}\equiv0$ \\
$\mathbf{E}_{6,11}$, $\mathbf{E}_{6,12}$, $\mathbf{E}_{6,22}$ & second gradients
  & the $\sigma_{33}$ interior build-up
  & $\sigma_{33}$ keeps only its face terms \\
$\mathbf{qt}_6$ & the \emph{applied} surface-pressure ladder
  \texttt{[q q,1 q,2 q,11 q,12 q,22]}
  & the face traction $\sigma_{33} = -q$ & $\sigma_{33}=0$ at that face \\
\bottomrule
\end{tabular}
\caption{The four driver groups of the through-thickness recovery.}
\label{tab:solid-drivers}
\end{table}

$Q_1$ and $Q_2$ are deliberately not part of $\mathbf{F}_6$: there is no shear
slot among the 42 drivers, because in the asymptotic construction the transverse
shears are produced by the gradient brackets. They enter only as the per-station
rescale targets $\int\sigma_{13}\,dz = Q_1$ and $\int\sigma_{23}\,dz = Q_2$, so
the profile \emph{shape} comes from the gradients while the carried force is
exactly the one you supplied.

\subsection{Driver A: one station, resultants typed into the script}

\texttt{rm\_dehom.py} recovers a single station. The input is the
\texttt{USER INPUT} block at the bottom of the file: \texttt{SG\_YAML} (the mesh
written in Section~\ref{sec:solid-plate1d}), \texttt{F6} and \texttt{Q2v}, plus
the optional gradient and load ladders. \texttt{F6} and \texttt{Q2v} ship as
\texttt{None} and the script raises a \texttt{ValueError} rather than run, by
design --- a built-in number would silently masquerade as a result. The block
below is the file with the two mandatory entries filled in with an example
loading; the values are yours to choose.

\begin{lstlisting}[style=opensg,caption={The USER INPUT block of rm\_dehom.py, with the two mandatory entries supplied.}]
SG_YAML = os.path.join(HERE, "1dsg.yaml")
F6  = [1.0e6, 0.0, 0.0, 5.0e4, 0.0, 0.0]  # [N11, N22, N12, M11, M22, M12]
Q2v = [0.0, 0.0]                          # [Q1, Q2]  [N/m]
DE1 = DE2 = None                          # d(E6)/dx, d(E6)/dy
D11 = D12 = D22 = None                    # second gradients
QT6 = None                                # [q q,1 q,2 q,11 q,12 q,22]
LATTICE = "gauss"                         # "gauss" (default) or "nodes"
\end{lstlisting}

\begin{lstlisting}[style=shellst]
cd examples/OpenSG-solid/2_get_plate_stress_from_1DSG
python rm_dehom.py
\end{lstlisting}

It writes \texttt{1dsg\_dehom.SM}, \texttt{.EM}, \texttt{.U} and \texttt{.vtk}
--- the VTK is part of the default output, not an extra step --- and prints
\texttt{max|S11|} through \texttt{max|S23|}. Because the script refuses to run
until you supply \texttt{F6} and \texttt{Q2v}, these four files are produced by
your own run and are not checked into the folder.

With every optional input left at \texttt{None} this is the honest classical
recovery: $\sigma_{11}, \sigma_{22}, \sigma_{12}$ exact layer by layer,
$\sigma_{13} = \sigma_{23} = 0$ exactly, $\sigma_{33}$ at machine zero. One
station has no neighbours, so it cannot supply gradients, and the code does not
invent them. \texttt{LATTICE} selects the output depths: \texttt{"gauss"} (the
default) puts them at the Gauss abscissae of each element, \texttt{"nodes"} at
the SG nodes.

\subsection{Driver B: a whole field from a \texttt{.ff} station table}

\texttt{rm\_dehom\_ff.py} is the general driver and the one to use for real
work. A \texttt{.ff} file is the plate analogue of a VABS \texttt{.glb}:
whitespace separated, \texttt{\#} comments, one row per in-plane station, in any
row order and on any station geometry. The shipped table was collected from an
Abaqus S8R shell field, but any solver can write it --- the driver contains no
Abaqus. Its first three lines state the format:

\begin{lstlisting}[style=shellst]
# plate dehom FF station table (collected from ex5_shell_S8R_field.dat)
# x[m] y[m] N11 N22 N12[N/m] M11 M22 M12[N] Q1 Q2[N/m] u1 u2 u3[m] wx wy[-]
# 40 x 40 rectangular station lattice, x-fastest
\end{lstlisting}

Ten columns are required (\texttt{x y N11 N22 N12 M11 M22 M12 Q1 Q2}); the
optional columns 11--15 add the mid-surface displacement $u_1, u_2, u_3$ and
slopes $w_x, w_y$, and without them the \texttt{.U} file carries the warping
alone. The loader accepts exactly 10 or 15 columns and raises otherwise.
\textbf{Every gradient driver is computed from the stations themselves}, by one
of two paths chosen automatically: exact tensor-grid finite differences
(\texttt{np.gradient} on the true non-uniform spacings) when the coordinates
form a complete rectangular lattice, or a meshfree KD-tree patch fit
($k = 12$ nearest neighbours, distance-weighted local quadratic least squares,
first and second derivatives from one batched solve) for any other cloud. The
flag \texttt{--scatter} forces the meshfree path.

Unlike the single-station driver, \texttt{rm\_dehom\_ff.py} has no
\texttt{LATTICE} variable at all: it always calls
\texttt{z\_stations(sg, "gauss")} and therefore always reports at the element
Gauss depths. \texttt{LATTICE = "nodes"} exists only in \texttt{rm\_dehom.py}.

The applied face pressure is \emph{not} in the table --- loads are input, not
solver output --- so it lives in \texttt{layup\_db.yaml}:

\begin{lstlisting}[style=opensg,caption={The dehom: block of layup\_db.yaml.}]
dehom:
  ff: Abaqus_results/ex5_shell_S8R_field.ff
  loads:
    top: {q0: 68.9476e6, shape: sinsin}   # uniform | sin_x | sin_y | sinsin
    bottom: none
\end{lstlisting}

Each \texttt{shape} fixes the whole six-term ladder analytically, differentiated
with \texttt{a} and \texttt{b} from the \texttt{plate:} block ($a = b = 1.524$~m
in the shipped file). Bottom-face loads raise \texttt{NotImplementedError}: the
42-driver operator block carries the top ladder only.

\begin{lstlisting}[style=shellst]
cd examples/OpenSG-solid/2_get_plate_stress_from_1DSG
python rm_dehom_ff.py ex5_shell_S8R_field.ff
\end{lstlisting}

Pass the table explicitly, as above. The \texttt{dehom: ff:} key in the shipped
\texttt{layup\_db.yaml} points at
\texttt{Abaqus\_results/ex5\_shell\_S8R\_field.ff}, and that folder does not
exist in this example: the table sits at the top level of the folder as
\texttt{ex5\_shell\_S8R\_field.ff}. Leaving the path to the key therefore trips
the driver's own guard, which reads \texttt{no .ff table: pass it on the command
line or set dehom: ff in} followed by the database path. Either give the file
name on the command line, as above, or correct the key. \texttt{--db} points at a database elsewhere. The run prints its chosen
gradient path (\texttt{lattice FD (40 x 40)} for this table), the operator
build, the \texttt{jax.vmap} over stations, and the peak of each stress
component.

\subsection{What comes out}

Outputs are named after the table, \texttt{<ff stem>\_ff}.

\begin{table}[htbp]
\centering
\footnotesize
\begin{tabular}{lp{8.0cm}}
\toprule
file & content \\
\midrule
\texttt{ex5\_shell\_S8R\_field\_ff.SM} & 3-D stress: columns \texttt{x y z} then
  \texttt{S11 S22 S33 S12 S13 S23} in Pa \\
\texttt{ex5\_shell\_S8R\_field\_ff.EM} & 3-D strain, same layout, \texttt{E11}
  through \texttt{E23} \\
\texttt{ex5\_shell\_S8R\_field\_ff.U} & 3-D displacement \texttt{U1 U2 U3} in m \\
\texttt{ex5\_shell\_S8R\_field\_ff\_gauss.vtk} & ASCII
  \texttt{STRUCTURED\_GRID}, \texttt{DIMENSIONS 40 40 45}, 72\,000 points,
  written only on a lattice \\
\bottomrule
\end{tabular}
\caption{Outputs of \texttt{rm\_dehom\_ff.py} on the shipped table.}
\label{tab:solid-ex2out}
\end{table}

The three text files share a two-line header naming the recovery and its station
set. For this run the second line reads \texttt{1600 .ff stations
[lattice FD (40 x 40)] x the 45 SG-element Gauss depths; ALL gradient drivers
from the stations; S/E in the PLY MATERIAL frame (U global); z = 0 is the
reference surface}. Four conventions matter when you read the files.

\begin{itemize}
\item $z$ is in the \textbf{SG frame}: $z = 0$ is the reference surface, not the
  mid-plane, unless \texttt{fraction} put them at the same place.
\item The depths are the Gauss abscissae of each five-noded element --- five per
  ply here, 45 in all --- so every output point is strictly inside one ply, and
  a material interface, where stress is genuinely double-valued, never appears
  in a file.
\item Stress and strain are written in the \textbf{ply material frame} of each
  depth, which is where failure criteria live. The chain computes in the
  laminate frame --- it must, since the gradients and the $Q$-rescale integrate
  across plies --- and rotates only at write time. The displacement \texttt{.U}
  stays global.
\item Component order in the files is \texttt{11 22 33 12 13 23}, reordered from
  the internal SG Voigt order $(11, 22, 33, 23, 13, 12)$ on write.
\end{itemize}

The \texttt{\_ff\_gauss.vtk} is the file to open when you want to look at the
result rather than read it. Figure~\ref{fig:solid-pv1dsg} is the shipped table
recovered and rendered straight out of it.

\begin{figure}[htbp]
\centering
\includegraphics[width=0.47\linewidth]{\FIGDIR/pv_plate_1dsg_s11.png}\hfill
\includegraphics[width=0.47\linewidth]{\FIGDIR/pv_plate_1dsg_s13.png}
\caption{ParaView renders of the recovered field on the SG, produced from the
shipped \texttt{ex5\_shell\_S8R\_field\_ff\_gauss.vtk} by
\texttt{docs/latex/render\_pv.py}. Left: the \texttt{S11} array, $\sigma_{11}$
in Pa, colour scale $\pm5.23\times10^{10}$; right: the \texttt{S13} array,
$\sigma_{13}$ in Pa, colour scale $\pm8.38\times10^{8}$. The grid points
\emph{are} the recovery stations --- the $40\times40$ in-plane lattice times the
45 SG-element Gauss depths --- so each coloured value is one recovered
Gauss-point value, drawn with flat shading and never averaged across an element
boundary. The camera looks down the thickness axis, so the visible face is the
outermost Gauss depth. $\sigma_{11}$ there follows the sin-sin pressure of the
\texttt{dehom: loads:} block and peaks at the plate centre; $\sigma_{13}$ is
near zero on the full-field scale, because the transverse shear must vanish at
the traction-free face, and changes sign across mid-span in $x$ following the
$\cos(\pi x/a)$ gradient of that same load.}
\label{fig:solid-pv1dsg}
\end{figure}

\subsection{What it is validated against}

The accuracy table in
\texttt{src/opensg\_solid/rm\_plate\_1D/USER\_GUIDE.md} records the recovery
errors against reference elasticity solutions, falling at $\mathcal{O}(S^{-2})$
in the slenderness $S$.

\begin{table}[htbp]
\centering
\small
\begin{tabular}{lccl}
\toprule
component (chain) & $S = 4$ & $S = 10$ & thin \\
\midrule
$\sigma_{13}$ (cross-ply) & 21\,\% & 4.4\,\% & 0.18\,\% at $S = 50$ \\
$\sigma_{33}$ equilibrium & 5.5\,\% & 1.2\,\% & 0.05\,\% \\
$\sigma_{11}$ second-order & 27\,\% & 2.0\,\% & 0.04\,\% at $S = 64$ \\
$\sigma_{12}$ angle-ply & 18\,\% & --- & 0.07\,\% at $S = 64$ \\
$U_1$ & 30\,\% & 1.4\,\% & 0.02\,\% at $S = 64$ \\
\bottomrule
\end{tabular}
\caption{Recovery accuracy against reference elasticity solutions, from
\texttt{USER\_GUIDE.md}.}
\label{tab:solid-accuracy}
\end{table}

Soft-core sandwiches need a larger $S$ for the in-plane second-order recovery;
their $\sigma_{13}$, $\sigma_{33}$ and displacements stay accurate at $S = 10$.
Separately, \texttt{dehom\_README.md} in the example folder records that the
finite-difference gradient path and the analytically differentiated known-mode
path were run side by side on the same field and agree to better than
$1\,\%$ on every component in the interior (four integration points trimmed per
edge); the $11$--$15\,\%$ whole-field figure is entirely the one-sided stencil
band about two elements wide at the boundary. That same README derives the full
chain equation by equation and states which term is Yu's and which is this
project's. The second-order recovery gate itself is
\texttt{examples/OpenSG-solid/CS\_OpeNSG\_exampels/static/validate\_v2.py},
which drives four Pagano cylindrical-bending cases with exact closed-form
harmonic drivers; its reference \texttt{.dat} profiles are produced by the
per-case \texttt{run\_case.py} and are not checked in.

\section{Plate properties from a 2-D cell SG}
\label{sec:solid-plate2d}

\subsection{What it computes}

When the wall has in-plane microstructure --- a honeycomb core, a corrugation, a
truss --- the gene is a 2-D mesh of one periodic cell, and \texttt{plate\_homo\_2d}
homogenizes it to the plate ABD with \texttt{n\_model=2}. Opposite faces, edges
and corners are tied through the sparse periodic map;
\texttt{boundary="periodic"} is the default and is the correct treatment for a
periodic medium.

Everything below the runner happens inside the single \texttt{plate\_homo\_2d}
call, in this order.

\captree{
  \node[inputnode] (i1) {\texttt{RHC\_SW\_2UC\_45.yaml}\\
    the 2-D cell SG (a SwiftComp\\ \texttt{.sc} is converted once by\\
    \texttt{helper.sc\_to\_yaml.convert})};
  \node[funcnode, below=of i1] (f1) {\texttt{sg\_homo.plate\_homo\_2d}\\
    $\rightarrow$ \texttt{sg\_mesh.load\_sg\_input}\\
    the parsed dict: \texttt{dim} (inferred\\ from the mesh),
    \texttt{nodes}, \texttt{cells},\\ \texttt{mat\_id}, \texttt{materials}};
  \node[funcnode, right=of f1] (p1) {\texttt{sg\_mesh.plot\_sg\_mesh}};
  \node[funcnode, below=of f1] (f2) {\texttt{fe\_jax.setup.}\\
    \texttt{mesh\_to\_periodic\_sparse\_assembly\_map}\\
    ties opposite faces, edges, corners};
  \node[outnode, right=of f2] (op) {\texttt{RHC\_SW\_2UC\_45\_mesh.png}};
  \node[funcnode, below=of f2] (f3)
    {\texttt{sg\_materials.get\_heterogeneous\_C\_matrix}\\
     \texttt{C\_ess} $(E,6,6)$, rebuilt from\\
     \texttt{material\_param} and \texttt{angles}};
  \node[inputnode, right=of f3] (i2) {\texttt{n\_model = 2}, \texttt{name},\\
    \texttt{material\_param}, \texttt{angles}\\ in \texttt{plate\_homo\_2dsg.py}};
  \node[funcnode, below=of f3] (f4) {\texttt{sg\_homo.\_homo\_direct}\\
    six unit macro states, one sparse\\
    factorization $\rightarrow$ \texttt{V0}, \texttt{C\_eff}, \texttt{omega}};
  \node[funcnode, right=of f4] (f5) {\texttt{sg\_homo.plate\_shear\_ladder}\\
    only at \texttt{shear\_refined=True}\\
    $\rightarrow$ \texttt{G\_msg}, \texttt{ABDG}};
  \node[datanode, below=of f4] (d1) {\texttt{r["C\_eff"]} $6\times6$ ABD\\
    \texttt{[N11 N22 N12 M11 M22 M12]}};
  \node[funcnode, below=of d1] (f6) {\texttt{sg\_homo.write\_sc\_K}};
  \node[outnode, right=of f6] (o1) {\texttt{RHC\_SW\_2UC\_45.out}};
  \draw[capedge] (i1) -- (f1);
  \draw[capedge] (f1) -- (p1);
  \draw[capedge] (p1) -- (op);
  \draw[capedge] (f1) -- (f2);
  \draw[capedge] (f2) -- (f3);
  \draw[capedge] (i2) -- (f3);
  \draw[capedge] (f3) -- (f4);
  \draw[capedge2] (f4) -- (f5);
  \draw[capedge] (f4) -- (d1);
  \draw[capedge] (d1) -- (f6);
  \draw[capedge] (f6) -- (o1);
}

\subsection{The example SG and the input you edit}

\texttt{RHC\_SW\_2UC\_45.yaml} is a two-unit-cell composite honeycomb sandwich
section: 4251 nodes, 6640 three-node triangles, three materials. Coordinates
are in millimetres and moduli in MPa. The yaml \emph{is} the whole problem: a
short scalar header above the mesh blocks says what to run, and the mesh and
material blocks say what to run it on.

\begin{lstlisting}[style=opensg,caption={The header of a solid SG yaml. Every key is defaulted; a headerless mesh still runs.}]
n_model: 2      # 1 = beam, 2 = plate, 3 = solid -- the macro model
refined: 0      # 0 = classical (plate ABD / beam EB); 1 = shear-refined
analysis: H     # H = homogenization | D = dehomogenization
\end{lstlisting}

\begin{table}[htbp]
\centering
\footnotesize
\begin{tabular}{llp{7.6cm}}
\toprule
key & values & meaning \\
\midrule
\texttt{msg} & \texttt{shell}, \texttt{solid}
  & which engine owns the file; omit and the mesh dialect decides \\
\texttt{n\_model} & \texttt{1}, \texttt{2}, \texttt{3}
  & macro model: 1 beam, 2 plate, 3 3-D solid \\
\texttt{refined} & \texttt{0}, \texttt{1}
  & 0 classical (plate ABD $6\times6$, beam EB $4\times4$), 1 shear-refined
    (plate ABDG $8\times8$, beam Timoshenko $6\times6$); ignored for the solid
    model \\
\texttt{analysis} & \texttt{H}, \texttt{D}
  & optional; the command-line argument wins \\
\texttt{epsilon\_bar} & 6 floats & the macro state for \texttt{D} \\
\texttt{omega} & float
  & optional user SG measure, overriding the one measured from the mesh \\
\texttt{aperiodic} & \texttt{1}
  & optional; zero fluctuation on every bounding-box-face node instead of
    periodicity \\
\bottomrule
\end{tabular}
\caption{The header contract of the solid SG yaml. \textbf{There is no
\texttt{dim:} key and no \texttt{scale:} key.} The SG dimension is read from
the mesh --- the element node count, with ties broken by the leading
node-coordinate columns that actually vary --- and the SG measure $\omega$ is
measured from the mesh. A legacy \texttt{dim:}/\texttt{scale:} left in an old
file is still tolerated by the loader, but neither is part of the input
contract and neither should be written into a new file. (The
\texttt{junction:} family of keys belongs to \texttt{msg-shell} recovery and is
not read on this route.)}
\label{tab:solid-header}
\end{table}

The \texttt{material\_param} and \texttt{angles} arguments of
\texttt{plate\_homo\_2d} are an \emph{override} of the material blocks stored
in the file, kept for scripting; examples 4 and 7 drive the same cell that way:

\begin{lstlisting}[style=opensg,caption={The material override block, as examples 4 and 7 carry it.}]
n_model = 2                    # 1: Beam; 2: Plate; 3: 3D elastic
name = "RHC_SW_2UC_45"         # reads <name>.yaml

material_param = jnp.array([
    (108e3, 8e3, 8e3, 4e3, 4e3, 3e3, 0.32, 0.32, 0.30),      # ply +45
    (108e3, 8e3, 8e3, 4e3, 4e3, 3e3, 0.32, 0.32, 0.30),      # ply -45
    (69e3, 69e3, 69e3, 26.54e3, 26.54e3, 26.54e3, 0.30, 0.30, 0.30)])
angles = jnp.array([45.0, -45.0, 0.0])   # deg per material; 0.0 = none
\end{lstlisting}

\texttt{material\_param} is one row of engineering constants per material,
$(E_1, E_2, E_3, G_{12}, G_{13}, G_{23}, \nu_{12}, \nu_{13}, \nu_{23})$, and
\texttt{angles} one rotation per material in degrees. Passing them
\emph{overrides} the material blocks stored in the SG file. The shipped
\texttt{RHC\_SW\_2UC\_45.yaml} keeps the same data in exactly one place
instead: all three materials are \texttt{type: 1}, the nine engineering
constants together with their ply rotation --- \texttt{angle: 45.0} and
\texttt{angle: -45.0} for the two face laminae, and the isotropic aluminium
core ($E = 69000$~MPa, $G = 26540$~MPa, $\nu = 0.3$) at \texttt{angle: 0.0}.
A file that stores its plies \emph{pre-rotated} instead, as a \texttt{type: 2}
full $6\times6$ $\mathbf{C}$ block, must not be given \texttt{angles} as well,
or the plies are rotated twice. The same numbers are recorded next to the
runner in \texttt{mat\_override.yaml} for reference; omit both arguments to use
the file's own blocks unchanged.

The solver reads \texttt{<name>.yaml}. A SwiftComp \texttt{.sc} is converted
once with the packaged helper. That helper is a library function, not a
command-line tool: the module \texttt{sc\_to\_yaml} has no \texttt{\_\_main\_\_}
block, so
\texttt{python -m opensg\_solid.helper.sc\_to\_yaml file.sc}
does nothing. Call
\texttt{convert} instead, which writes \texttt{<base>.yaml} and
\texttt{<base>.msh} and returns the parsed dictionary:

\begin{lstlisting}[style=opensg,caption={Converting a SwiftComp .sc once.}]
from opensg_solid.helper.sc_to_yaml import convert

sc = convert("RHC_SW_2UC_45.sc")   # writes RHC_SW_2UC_45.yaml + .msh
\end{lstlisting}

\subsection{Run it}

The yaml says what to do, so the run is one command with one argument:

\begin{lstlisting}[style=shellst]
cd examples/OpenSG-solid/3_get_plate_props_from_2D_SG
opensg RHC_SW_2UC_45.yaml
\end{lstlisting}

\texttt{opensg} is the unified entry point: it resolves the engine from the
yaml's \texttt{msg:} key --- or, absent that key, from the mesh dialect --- and
hands the arguments on unchanged. \texttt{opensg\_solid RHC\_SW\_2UC\_45.yaml}
and \texttt{python -m opensg\_solid RHC\_SW\_2UC\_45.yaml} are the same run
through the engine-specific entry point.
\texttt{plate\_homo\_2dsg.py} in the folder is the identical call through the
Python API, \texttt{plate\_homo\_2d(name + ".yaml")}, for driving it from your
own script.

\subsection{What comes out}

The run writes \texttt{RHC\_SW\_2UC\_45.out}, the timed \texttt{Classical Plate}
stiffness and compliance in the \texttt{.K} layout, and
\texttt{RHC\_SW\_2UC\_45\_mesh.png}, the cell mesh with elements coloured by
material. It also prints the same $6\times6$ to the terminal, under the macro
law's own title. The \texttt{.out} banner records the SG measure and the
boundary treatment:

\begin{lstlisting}[style=shellst,basicstyle=\ttfamily\tiny]
 OpenSG msg-solid plate model, omega 35.248001, periodic

 The Effective Classical Plate Stiffness Matrix
 --------------------------------------------
     3.3585973E+005     7.7500043E+004    -1.4366892E+002     4.1488404E-001     1.4861902E+000     3.6602858E-002
\end{lstlisting}

closing with \texttt{ Time taken: 7.83 sec} for this 4251-node, 6640-triangle
cell.

\subsection{What the numbers mean}

$\omega = 35.248001$ is the cell's in-plane period, and it is directly checkable
against the mesh: the nodes run from $y_1 = -17.624$ to $+17.624$~mm. It is
\emph{measured}, never declared --- the yaml carries no normalization key,
because the homogenizer always measures its own SG. For a cell converted from
SwiftComp the measurement is also checkable against the source: the trailing
line of \texttt{RHC\_SW\_2UC\_45.sc} reads \texttt{35.248}, and the measured
$\omega$ agrees with it to the digits that line carries, which is the check
that the mesh and the \texttt{.sc} header describe the same cell. The
diagonal, in the mesh's own units:

\begin{table}[htbp]
\centering
\small
\begin{tabular}{lll}
\toprule
term & value & units \\
\midrule
$A_{11}$ & $3.3585973\times10^{5}$ & N/mm \\
$A_{22}$ & $1.1291139\times10^{5}$ & N/mm \\
$A_{66}$ & $1.0551654\times10^{5}$ & N/mm \\
$D_{11}$ & $9.0367521\times10^{7}$ & N$\cdot$mm \\
$D_{22}$ & $5.7011050\times10^{7}$ & N$\cdot$mm \\
$D_{66}$ & $4.6026320\times10^{7}$ & N$\cdot$mm \\
\bottomrule
\end{tabular}
\caption{Diagonal of the honeycomb cell's plate ABD.}
\label{tab:solid-ex3diag}
\end{table}

The membrane law is strongly anisotropic, $A_{11}/A_{22} = 2.97$, while the
bending law is much closer to isotropic, $D_{11}/D_{22} = 1.59$: the two
directions see the cell very differently in stretching and much less
differently in bending. The largest entry of the whole $B$ block is $4.954$
against $A$ terms of order $10^{5}$ and $D$ terms of order $10^{7}$, which is
what a cell symmetric about its mid-plane should give and is the first health
check to make on a run of your own.

\begin{figure}[htbp]
\centering
\includegraphics[width=0.78\textwidth]{\FIGDIR/rhc_2dsg_mesh.png}
\caption{The two-unit-cell honeycomb sandwich structure gene, elements coloured
by material, as written to \texttt{RHC\_SW\_2UC\_45\_mesh.png}. The
$\pm45^\circ$ face laminae (materials 1 and 2) are the thin horizontal bands at
$y_2 = \pm23$~mm and the isotropic core web (material 3) folds between them over
the in-plane period. The through-thickness coordinate is $y_2$, spanning
$\pm23.35$~mm.}
\label{fig:solid-rhcmesh}
\end{figure}

\subsection{The Reissner-Mindlin option}

Setting \texttt{refined: 1} in the yaml header --- or passing
\texttt{refined=1} to \texttt{plate\_homo\_2d}, of which
\texttt{shear\_refined=True} is the legacy plate-only spelling --- additionally
runs the RM first-order warping ladder and returns \texttt{r["G\_msg"]} (the
$2\times2$ transverse-shear block), \texttt{r["ABDG"]} (the same $8\times8$
block form as Eq.~\eqref{eq:solid-abdg}), \texttt{r["A6\_ladder"]} and
\texttt{r["Ustar\_rel"]}. When the ladder produces an SPD fit the \texttt{.out}
is written from the $8\times8$ and its matrices are titled
\texttt{Reissner-Mindlin} instead of \texttt{Classical Plate}; otherwise
the run falls back to the $6\times6$. The same header switch serves the beam
route: \texttt{n\_model: 1} with \texttt{refined: 0} reports the classical
\texttt{Euler-Bernoulli Beam} $4\times4$, and with \texttt{refined: 1} the
\texttt{Timoshenko Beam} $6\times6$. The solid law has no refined variant.

\begin{lstlisting}[style=opensg,caption={Requesting the 8x8 from a 2-D cell SG.}]
from opensg_solid.sg_homo import plate_homo_2d

r = plate_homo_2d("RHC_SW_2UC_45.yaml", material_param=material_param,
                  angles=angles, n_model=2, shear_refined=True)
print(r["ABDG"])
\end{lstlisting}

\subsection{The 1-D gene in the same folder, and a naming trap}

The folder also ships \texttt{Plate\_1D\_SG\_2UC\_45.yaml}, a \textbf{1-D}
structure gene --- two-node line elements through a stack, its dimension read
from the mesh rather than declared in the header --- which runs through the
identical command by naming that file instead.
\texttt{plate\_homo\_2d} reads 1-D, 2-D and 3-D genes from the same code path,
so a stack of plies can be homogenized either here or through the dedicated
through-thickness route of Section~\ref{sec:solid-plate1d}.

Be careful about the output name. \texttt{plate\_homo\_2d} always derives its
report path from the \emph{input} basename, so setting
\texttt{name = "Plate\_1D\_SG\_2UC\_45"} writes
\texttt{Plate\_1D\_SG\_2UC\_45.out} --- the \texttt{.K}-layout report --- and
nothing else. The file \texttt{Plate\_1D\_SG\_2UC\_45\_plate\_ABD.out} sitting
in that folder is not produced by \texttt{plate\_homo\_2dsg.py} at all: it is an
\texttt{np.savetxt} leftover from an older script style, a bare $6\times6$ under
a single comment line (\texttt{\# plate 6x6 [N11 N22 N12 M11 M22 M12] from the
1D SG Plate\_1D\_SG\_2UC\_45.sc}) with no banner, no compliance block and no
timing. Its contents are still a valid record of that gene ---
$A_{11} = A_{22} = 1.65452926\times10^{6}$,
$A_{66} = 6.42527474\times10^{5}$,
$D_{11} = D_{22} = 7.86727254\times10^{7}$,
$D_{66} = 3.12496023\times10^{7}$ in N/mm and N$\cdot$mm --- but do not expect
your run to reproduce that file name or that layout.

\section{Plate recovery from a 2-D SG}
\label{sec:solid-platedehom2d}

\subsection{What it computes}

Given a converged plate solution you know the macro plate strain at the point of
interest. Feeding it back into the homogenized cell recovers the local 3-D
strain and stress at \textbf{every Gauss point of the gene}, which is where a
honeycomb actually fails --- in the web, not in the smeared plate.
\texttt{dehom\_fields} returns the strain, the stress and the fluctuation
displacement in one pass, using the warping columns $V_0$ that the
homogenization already solved. Nothing is re-solved.

That reuse is the whole point of the tree below: the homogenization node runs
once and its result dictionary carries the warping modes, the element stiffness
table and the quadrature data straight into the recovery.

\captree{
  \node[inputnode] (i1) {\texttt{RHC\_SW\_2UC\_45.yaml}\\ the 2-D cell SG};
  \node[funcnode, below=of i1] (f1) {\texttt{sg\_homo.plate\_homo\_2d}\\
    \texttt{n\_model=2}, the Section~\ref{sec:solid-plate2d} solve};
  \node[funcnode, right=of f1] (f2) {\texttt{sg\_homo.write\_sc\_K},\\
    \texttt{sg\_mesh.plot\_sg\_mesh},\\
    \texttt{np.savetxt} in the runner};
  \node[datanode, below=of f1] (d1) {\texttt{r["C\_eff"]}, \texttt{r["V0"]},\\
    \texttt{C\_ess}, \texttt{x\_end}, \texttt{phi\_qn},\\
    \texttt{dphi\_dxi\_qnp}, \texttt{periodic\_cells\_en}};
  \node[outnode, right=of d1] (o1) {\texttt{RHC\_SW\_2UC\_45.out}\\
    \texttt{RHC\_SW\_2UC\_45\_mesh.png}\\
    \texttt{RHC\_SW\_2UC\_45\_plate\_ABD.out}};
  \node[funcnode, below=of d1] (f3) {\texttt{sg\_dehom.dehom\_fields}\\
    \texttt{\_dehom\_batch} and \texttt{\_u\_at\_gauss}\\
    $\rightarrow$ \texttt{Gam}, \texttt{Sig} $(E,Q,6)$, \texttt{U} $(E,Q,3)$};
  \node[inputnode, right=of f3] (i2) {\texttt{epsilon\_bar}\\
    \texttt{[e11 e22 2e12 k11 k22 2k12]}\\ in \texttt{plate\_dehom\_2dsg.py}};
  \node[funcnode, below=of f3] (f4) {\texttt{sg\_dehom.export\_gauss}\\
    \texttt{sg\_dehom.gauss\_coords},
    \texttt{\_write\_field}};
  \node[outnode, below=of f4] (o2) {\texttt{RHC\_SW\_2UC\_45\_dehom.txt}\\
    \texttt{.vtk}, \texttt{.SM}, \texttt{.EM}, \texttt{.U}};
  \draw[capedge] (i1) -- (f1);
  \draw[capedge] (f1) -- (f2);
  \draw[capedge] (f2) -- (o1);
  \draw[capedge] (f1) -- (d1);
  \draw[capedge] (d1) -- (f3);
  \draw[capedge] (i2) -- (f3);
  \draw[capedge] (f3) -- (f4);
  \draw[capedge] (f4) -- (o2);
}

\subsection{The input you edit}

The recovery needs one thing beyond the homogenization: the macro plate strain,
in the order \texttt{[e11, e22, 2e12, k11, k22, 2k12]}. On the command-line
route it is a header key of the SG yaml, read when the analysis is \texttt{D}:

\begin{lstlisting}[style=opensg,caption={The macro plate strain driving the recovery, in the yaml header.}]
epsilon_bar: [0.0, 0.1, 0.0, 0.1, 0.0, 0.0]
\end{lstlisting}

\texttt{plate\_dehom\_2dsg.py} sets the same six numbers in code, as
\texttt{epsilon\_bar = jnp.array([0.0, 0.1, 0.0, 0.1, 0.0, 0.0])}, alongside
the \texttt{name}, \texttt{material\_param} and \texttt{angles} of
Section~\ref{sec:solid-plate2d}.

The shipped state combines a transverse extension $\varepsilon_{22} = 0.1$ with
a bending curvature $\kappa_{11} = 0.1$~mm$^{-1}$. This vector is yours: it
comes from your plate solution at the station you care about, and there is no
default that would be meaningful.

\subsection{Run it}

The trailing \texttt{D} asks for dehomogenization: the run homogenizes first,
then recovers, and reports one \texttt{Time taken} covering the whole thing.

\begin{lstlisting}[style=shellst]
cd examples/OpenSG-solid/4_get_plate_dehom_from_2DSG
opensg RHC_SW_2UC_45.yaml D
\end{lstlisting}

\texttt{python plate\_dehom\_2dsg.py} is the same chain through the Python API,
and additionally writes the bare $6\times6$
\texttt{RHC\_SW\_2UC\_45\_plate\_ABD.out} through \texttt{np.savetxt}.

\subsection{What comes out}

\begin{table}[htbp]
\centering
\footnotesize
\begin{tabular}{lp{8.2cm}}
\toprule
file & content \\
\midrule
\texttt{RHC\_SW\_2UC\_45.out} & the timed \texttt{.K}-layout
  \texttt{Classical Plate} report written by \texttt{plate\_homo\_2d} inside the
  script; regenerated on each run \\
\texttt{RHC\_SW\_2UC\_45\_plate\_ABD.out} & the plain-text $6\times6$ ABD via
  \texttt{np.savetxt}, header \texttt{plate 6x6 [N11 N22 N12 M11 M22 M12] from
  the 2D SG RHC\_SW\_2UC\_45} \\
\texttt{RHC\_SW\_2UC\_45\_dehom.txt} & one row per Gauss point:
  \texttt{GP\_Index}, \texttt{Coord\_X}, \texttt{Coord\_Y}, then
  \texttt{Eps\_xx Eps\_yy Eps\_zz Eps\_yz Eps\_xz Eps\_xy} and \texttt{Sig\_xx}
  through \texttt{Sig\_xy} --- SwiftComp component order; one header line plus
  19\,920 rows \\
\texttt{RHC\_SW\_2UC\_45\_dehom.vtk} & ASCII \texttt{UNSTRUCTURED\_GRID} point
  cloud of the same 19\,920 points, each strain and stress component as its own
  \texttt{SCALARS} array, ready for ParaView \\
\texttt{RHC\_SW\_2UC\_45\_dehom.SM} / \texttt{.EM} / \texttt{.U} & the same
  fields in the \texttt{rm\_plate\_1D} layout: \texttt{x y z} then
  \texttt{S11 S22 S33 S12 S13 S23}, \texttt{E11} through \texttt{E23}, and
  \texttt{U1 U2 U3} \\
\texttt{RHC\_SW\_2UC\_45\_mesh.png} & the mesh figure, from the
  \texttt{plate\_homo\_2d} call the script makes first \\
\bottomrule
\end{tabular}
\caption{Outputs of \texttt{plate\_dehom\_2dsg.py}.}
\label{tab:solid-ex4out}
\end{table}

The 19\,920 rows are 6640 triangles times three quadrature points. The script
prints the ABD diagonal and the peak of each stress component.

Two conventions are worth knowing before reading the files. In the
\texttt{.SM}/\texttt{.EM}/\texttt{.U} tables the SG coordinates are placed in
the \textbf{last} slots of the \texttt{x y z} triple, so a 2-D gene leaves the
\texttt{x} column at zero and puts $y_1$ in \texttt{y} and the
through-thickness $y_2$ in \texttt{z}. And \texttt{.U} is the
\textbf{fluctuation-only} displacement --- its own header says so --- because a
unit-state SG recovery has no macro displacement to add.

\subsection{Reading the result}

The first Gauss point of the shipped run sits at
$(y_1, y_2) = (17.54067,\ 23.29375)$~mm, right at the top face of the cell whose
half-thickness is $23.35$~mm, and reports
%
\begin{equation}
\varepsilon_{11} = 2.329375, \qquad
\sigma_{11} = 8.755003\times10^{4}, \quad
\sigma_{22} = 7.102971\times10^{4}, \quad
\sigma_{12} = -6.542276\times10^{4}
\end{equation}
%
in MPa. That $\varepsilon_{11}$ is exactly $\kappa_{11} y_2
= 0.1 \times 23.29375$, the macro part of the strain reproduced to every printed
digit, which is the cheapest possible check that the strain state you passed is
the strain state that arrived. The other five components are what the recovery
adds: $\varepsilon_{22} = 0.2643$ and $\varepsilon_{33} = -0.9187$ are the
fluctuation response of the cell, and $\sigma_{12}$ is large even though the
applied macro shear strain $2\varepsilon_{12}$ is zero, which is consistent with
the off-axis face laminae the mesh figure shows at that height.

Over the whole gene the field maxima are:

\begin{table}[htbp]
\centering
\small
\begin{tabular}{lcc}
\toprule
component & $\max\lvert\sigma\rvert$ (MPa) & $\max\lvert\varepsilon\rvert$ \\
\midrule
$xx$ & $1.641632\times10^{5}$ & $2.329375$ \\
$yy$ & $9.199997\times10^{4}$ & $1.499038$ \\
$zz$ & $2.535821\times10^{4}$ & $1.032184$ \\
$yz$ & $1.083274\times10^{4}$ & $1.153314$ \\
$xz$ & $5.210321\times10^{3}$ & $8.807437\times10^{-1}$ \\
$xy$ & $7.369978\times10^{4}$ & $7.881358\times10^{-1}$ \\
\bottomrule
\end{tabular}
\caption{Component maxima over the 19\,920 Gauss points of the shipped plate
recovery.}
\label{tab:solid-ex4max}
\end{table}

The $xx$ strain maximum is again $\kappa_{11} y_2^{\max} = 2.329375$, as it must
be at the extreme fibre. Local states like the $\sigma_{zz}$ and $\sigma_{yz}$
columns are exactly what a smeared plate model cannot show you, and the reason
to run a dehomogenization at all. The driver here is a demonstration magnitude,
not a service load; scale it to your own macro state before reading the stresses
as failure indices.

\begin{figure}[htbp]
\centering
\includegraphics[width=0.62\linewidth]{\FIGDIR/pv_plate_dehom_sxx.png}
\caption{ParaView render of the recovered field on the SG mesh, produced from
the \texttt{RHC\_SW\_2UC\_45\_dehom.vtk} written by this example, using
\texttt{docs/latex/render\_pv.py}. The array is \texttt{Sig\_xx}, the axial
stress $\sigma_{xx}$ in MPa (units follow the SG input, millimetres and
megapascals), on a symmetric colour scale of $\pm1.64\times10^{5}$~MPa --- the
$xx$ row of Table~\ref{tab:solid-ex4max}. Every one of the 19\,920 points is a
single Gauss point carried by its own vertex cell and coloured flat, so no value
has been averaged across an element boundary. The field varies linearly with the
through-thickness coordinate $y_2$, which is the macro part $\kappa_{11}y_2$ of
the driver, and each face pack splits into two sub-bands of visibly different
$\sigma_{xx}$ at almost the same height, consistent with the $+45^\circ$ and
$-45^\circ$ laminae differing in the sign of their $C_{16}$ coupling.}
\label{fig:solid-pvplatedehom}
\end{figure}

\section{Equivalent 3-D solid from a 2-D SG}
\label{sec:solid-3dfrom2d}

\subsection{What it computes}

\texttt{n\_model=3} reduces an SG to an equivalent \textbf{3-D solid} stiffness:
the $6\times6$ material law of the homogenized continuum, in Voigt order
$[11, 22, 33, 23, 13, 12]$ with engineering shears. On a \textbf{2-D} SG this is
the generalized-plane-strain problem: the cell is periodic in its two in-plane
directions and uniform out of plane, six unit macro strains are applied, the
fluctuation field $\mathbf{V}_0$ is solved for each, and the effective law is
%
\begin{equation}
\mathbf{D}_{\rm eff} =
 \frac{\mathbf{D}_{ee} + \mathbf{V}_0^{\mathsf T}\mathbf{D}_{he}}{\omega},
\label{eq:solid-deff}
\end{equation}
%
with $\omega$ the cell area. This is the classical micromechanics unit-cell
problem, which makes it the natural place to validate the route against
published numbers.

The shipped runner takes the self-contained path on the left of the tree; the
engine path on the right is the same macro model reached through
\texttt{plate\_homo\_2d}, and is what a 2-D SG in the ordinary
\texttt{sg\_mesh.load\_sg\_input} dialect goes through.

\captree{
  \node[inputnode] (i1) {\texttt{UDcomp\_2D.msh}\\ gmsh 2.2, physical tags\\
    naming the matrix and fibre};
  \node[funcnode, below=of i1] (f1) {\texttt{make\_udcomp\_yaml.py}\\
    nodes, 1-based triangles, a constant\\
    \texttt{elementOrientations} frame, the two\\
    materials in MPa, \texttt{sets.element}};
  \node[outnode, below=of f1] (o1) {\texttt{UDcomp\_2D.yaml}};
  \node[funcnode, below=of o1] (f2) {\texttt{solid\_homo\_2dsg.C\_from\_eng}\\
    then the CST assembly:
    \texttt{Ke}, \texttt{Fe},\\ \texttt{Dee}, \texttt{Dhe}, \texttt{omega}};
  \node[inputnode, right=of f2] (i2) {\texttt{n\_model = 3}, \texttt{name},\\
    the \texttt{paper} dictionary\\ in \texttt{solid\_homo\_2dsg.py}};
  \node[funcnode, below=of f2] (f3) {\texttt{solid\_homo\_2dsg.pair}\\
    right to left, top to bottom;\\
    corner chains by \texttt{master = master[master]};\\
    three rigid translations pinned};
  \node[datanode, below=of f3] (d1) {\texttt{V0} $(n_{\rm dof},6)$\\
    $\mathbf{D}_{\rm eff} = (\mathbf{D}_{ee}
      + \mathbf{V}_0^{\mathsf T}\mathbf{D}_{he})/\omega$};
  \node[funcnode, right=of d1] (e1) {\texttt{sg\_homo.plate\_homo\_2d}\\
    \texttt{n\_model=3} on
    \texttt{RHC\_SW\_2UC\_45.yaml}\\ $\rightarrow$
    \texttt{RHC\_SW\_2UC\_45.out}\\ with the engineering-constants block};
  \node[outnode, below=of d1] (o2) {\texttt{UDcomp\_2D\_Deff.out}\\
    \texttt{UDcomp\_2D\_constants\_vs\_paper.dat}\\
    \texttt{UDcomp\_2D\_mesh.png}};
  \draw[capedge] (i1) -- (f1);
  \draw[capedge] (f1) -- (o1);
  \draw[capedge] (o1) -- (f2);
  \draw[capedge] (i2) -- (f2);
  \draw[capedge] (f2) -- (f3);
  \draw[capedge] (f3) -- (d1);
  \draw[capedge2] (d1) -- (e1);
  \draw[capedge] (d1) -- (o2);
}

\subsection{The example SG}

\texttt{examples/OpenSG-solid/5\_get\_solid\_props\_from\_SG} homogenizes a
unidirectional fibre/matrix square unit cell: 801 nodes, 1536 three-node
triangles, 512 matrix elements and 1024 fibre elements, cell area
$\omega = 1$.

\begin{table}[htbp]
\centering
\small
\begin{tabular}{lllll}
\toprule
material & $E_1, E_2, E_3$ (MPa) & $G_{12}, G_{13}, G_{23}$ (MPa)
 & $\nu_{12}, \nu_{13}, \nu_{23}$ & density \\
\midrule
matrix & 4760, 4760, 4760 & 1737.226 each & 0.37, 0.37, 0.37 & 1200.0 \\
fibre & 276000, 19500, 19500 & 70000, 70000, 5735 & 0.28, 0.28, 0.70 & 1800.0 \\
\bottomrule
\end{tabular}
\caption{The two phases of the UD unit cell. The matrix shear modulus is written
by the converter as $4760/(2\cdot1.37)$.}
\label{tab:solid-udmat}
\end{table}

\begin{figure}[htbp]
\centering
\includegraphics[width=0.5\textwidth]{\FIGDIR/udcomp_2dsg_mesh.png}
\caption{The unidirectional composite unit cell written to
\texttt{UDcomp\_2D\_mesh.png}, fibre and matrix elements coloured by phase, with
the SG coordinates labelled $y_2$ and $y_3$.}
\label{fig:solid-udmesh}
\end{figure}

\subsection{Step 1: build the SG yaml from the gmsh mesh}

\texttt{make\_udcomp\_yaml.py} is a one-time converter from
\texttt{UDcomp\_2D.msh} (gmsh 2.2 format) to \texttt{UDcomp\_2D.yaml}. It writes
the nodes as \texttt{- [x y z]} rows, the triangles as 1-based
\texttt{- [n1 n2 n3]} rows, a constant \texttt{elementOrientations} frame
($e_1$ = axial $z$, $e_2$ = $+x$, $e_3$ = $+y$), the two materials as
engineering constants in MPa, and \texttt{sets.element} splitting the mesh into
the matrix and fibre phases by gmsh physical tag. The yaml is already checked
in, so this step is only needed if you remesh.

\begin{lstlisting}[style=shellst]
cd examples/OpenSG-solid/5_get_solid_props_from_SG
python make_udcomp_yaml.py
\end{lstlisting}

\subsection{Step 2: homogenize}

\begin{lstlisting}[style=shellst]
python solid_homo_2dsg.py
\end{lstlisting}

The script is kept rather than replaced by the unified command because it does
more than homogenize: it writes the comparison table against the published
constants and the mesh figure. For a 2-D SG in the ordinary solid dialect the
same macro law is one command, \texttt{opensg <name>.yaml} with
\texttt{n\_model: 3} in the header.

The \texttt{User Input} block is three lines: \texttt{n\_model = 3},
\texttt{name = "UDcomp\_2D"} (which reads \texttt{<name>.yaml}), and the
\texttt{paper} dictionary of reference constants to compare against.

This particular runner is deliberately self-contained. It assembles
constant-strain triangles, the periodic master-slave node map (right to left
matching $y$, top to bottom matching $x$, with corner chains resolved by three
passes of \texttt{master = master[master]}) and the three pinned rigid
translations in plain NumPy, so the entire \texttt{n\_model=3} algebra is
readable in one file. It therefore writes its own \texttt{\_Deff.out} through
\texttt{np.savetxt} rather than the SwiftComp-layout report, and it consumes the
\texttt{nodes}/\texttt{elements}/\texttt{elementOrientations}/\texttt{materials}/\texttt{sets}
yaml dialect that \texttt{make\_udcomp\_yaml.py} emits, not the
\texttt{nodes}/\texttt{cells}/\texttt{mat\_id}/\texttt{materials}
dialect that \texttt{opensg\_solid.sg\_mesh.load\_sg\_input} parses. The engine
path itself is exercised on the honeycomb SG: \texttt{n\_model=3} on
\texttt{RHC\_SW\_2UC\_45.yaml} runs through \texttt{plate\_homo\_2d} and writes
the full SwiftComp \texttt{.out} with its engineering-constants block.

\subsection{What comes out}

\begin{table}[htbp]
\centering
\footnotesize
\begin{tabular}{lp{8.2cm}}
\toprule
file & content \\
\midrule
\texttt{UDcomp\_2D\_Deff.out} & the $6\times6$ $\mathbf{D}_{\rm eff}$ in MPa,
  Voigt order \texttt{[e11 e22 e33 2e23 2e13 2e12]}, with $\omega$ in the
  header \\
\texttt{UDcomp\_2D\_constants\_vs\_paper.dat} & the nine effective engineering
  constants next to the published values and the percentage difference \\
\texttt{UDcomp\_2D\_mesh.png} & the unit cell, elements coloured by phase \\
\bottomrule
\end{tabular}
\caption{Outputs of \texttt{solid\_homo\_2dsg.py}.}
\label{tab:solid-ex5out}
\end{table}

\subsection{What the numbers mean, and what they were validated against}

The nine constants are read off the compliance
$\mathbf{S} = \mathbf{D}_{\rm eff}^{-1}$ as $E_i = 1/S_{ii}$,
$G_{ij} = 1/S_{kk}$ and $\nu_{ij} = -S_{ij}/S_{ii}$, and compared against Table
II Model-3 of the ASC-23 OpenSG-Solid paper --- the same case as the FEniCS
\texttt{3DModel.ipynb} in \texttt{wenbinyugroup/OpenSG}. The shipped
\texttt{UDcomp\_2D\_constants\_vs\_paper.dat} reads:

\begin{table}[htbp]
\centering
\small
\begin{tabular}{lrrr}
\toprule
constant & this code & paper & \% difference \\
\midrule
$E_1$ (GPa) & 167.255337 & 167.255300 & $+0.0000$ \\
$E_2$ (GPa) & 11.413657 & 11.413600 & $+0.0005$ \\
$E_3$ (GPa) & 11.413657 & 11.413600 & $+0.0005$ \\
$G_{23}$ (GPa) & 3.095555 & 3.095500 & $+0.0018$ \\
$G_{13}$ (GPa) & 6.801763 & 6.801700 & $+0.0009$ \\
$G_{12}$ (GPa) & 6.801763 & 6.801700 & $+0.0009$ \\
$\nu_{12}$ & 0.311780 & 0.311780 & $-0.0000$ \\
$\nu_{13}$ & 0.311780 & 0.311780 & $-0.0000$ \\
$\nu_{23}$ & 0.575743 & 0.575740 & $+0.0006$ \\
\bottomrule
\end{tabular}
\caption{\texttt{UDcomp\_2D\_constants\_vs\_paper.dat}, verbatim.}
\label{tab:solid-udvspaper}
\end{table}

Every constant agrees with the published value to within $0.002\,\%$, which is
the round-off of the paper's own printed precision. The route is exact against
its reference, not merely close.

The stiffness itself, from \texttt{UDcomp\_2D\_Deff.out} (MPa, $\omega = 1$):
%
\widematrix{
\mathbf{D}_{\rm eff} =
\begin{bmatrix}
1.726544\times10^{5} & 8.658488\times10^{3} & 8.658488\times10^{3} & 0 & 0 & 0\\
8.658488\times10^{3} & 1.750725\times10^{4} & 1.026390\times10^{4} & 0 & 0 & 0\\
8.658488\times10^{3} & 1.026390\times10^{4} & 1.750725\times10^{4} & 0 & 0 & 0\\
0 & 0 & 0 & 3.095555\times10^{3} & 0 & 0\\
0 & 0 & 0 & 0 & 6.801763\times10^{3} & 0\\
0 & 0 & 0 & 0 & 0 & 6.801763\times10^{3}
\end{bmatrix}}

Two health checks are worth making a habit of. First, the symmetry the geometry
demands is recovered: $D_{22} = D_{33}$ and $D_{55} = D_{66}$ to all printed
digits, and every normal-shear coupling term in the shipped file is either
exactly zero or below $3.4\times10^{-2}$~MPa, six to seven orders under the
diagonal, which is the numerical signature of a correctly tied periodic cell.
Second, the transverse plane is \textbf{not} isotropic: a transversely isotropic
material would satisfy $D_{44} = (D_{22} - D_{23})/2 = 3621.7$~MPa, whereas the
computed $D_{44}$ is $3095.6$~MPa, $14.5\,\%$ lower. A square fibre array is
tetragonal, so $G_{23}$ genuinely has to be homogenized rather than inferred
from $E_2$ and $\nu_{23}$, which is precisely why a unit-cell solve is needed
instead of a rule of mixtures.

\section{Equivalent 3-D solid from a 3-D SG}
\label{sec:solid-3dfrom3d}

\subsection{What it computes}

When the SG is periodic in all three directions --- a triply-periodic minimal
surface (TPMS) unit cell, a lattice block, any genuine 3-D microstructure ---
the same macro model is reached with
\texttt{plate\_homo\_2d(..., n\_model=3, boundary="periodic")} on a tet or hex
mesh. \texttt{"periodic"} is the default on every route;
\texttt{boundary="aperiodic"} maps the boundary solution onto the
bounding-box-face nodes as Dirichlet data (solid elements carry only
translations, so this clamps every degree of freedom of those nodes) and is the
explicit request for SGs that genuinely are not periodic. Aperiodic supports the
plate and solid models with \texttt{solver="direct"} only; asking for it with
\texttt{n\_model=1} or \texttt{solver="cg"} raises.

The tree is the sample runner and the boundary comparison, the two paths every
number in this section comes from. Note the loader branch: a SwiftComp
\texttt{.sc} is converted on the way in, so the \texttt{.yaml} and \texttt{.msh}
appear in the working directory as a side effect of the homogenization.

\captree{
  \node[inputnode] (i1) {\texttt{Sample\_1.sc} or \texttt{Sample\_2.sc}\\
    SwiftComp linear tets on a unit cell};
  \node[inputnode, right=of i1] (i2) {\texttt{n\_model = 3},
    \texttt{E}, \texttt{nu},\\ \texttt{boundary},
    \texttt{workdir} in\\ \texttt{run\_sample\_N.py}};
  \node[funcnode, below=of i1] (f1) {\texttt{sg\_homo.plate\_homo\_2d}\\
    the single entry point};
  \node[funcnode, below=of f1] (f2) {\texttt{sg\_mesh.load\_sg\_input}\\
    $\rightarrow$ \texttt{helper.sc\_to\_yaml.convert}};
  \node[funcnode, right=of f2] (p1) {\texttt{sg\_mesh.plot\_sg\_mesh}};
  \node[funcnode, below=of f2] (f3) {\texttt{fe\_jax.setup.}\\
    \texttt{mesh\_to\_periodic\_sparse\_assembly\_map},\\
    or the bounding-box-face Dirichlet set\\
    at \texttt{boundary="aperiodic"}; then\\
    \texttt{sg\_materials.get\_heterogeneous\_C\_matrix}};
  \node[outnode, right=of f3] (oa) {\texttt{Sample\_N.yaml}\\
    \texttt{Sample\_N.msh}\\ \texttt{Sample\_N\_mesh.png}};
  \node[funcnode, below=of f3] (f4) {\texttt{sg\_homo.\_homo\_direct}\\
    six unit macro states $\rightarrow$\\
    \texttt{C\_eff}, \texttt{V0}, \texttt{omega}};
  \node[funcnode, right=of f4] (fb) {\texttt{compare\_boundary\_modes.py}\\
    both samples under both treatments};
  \node[funcnode, below=of f4] (f5) {\texttt{sg\_homo.write\_sc\_K}, then the\\
    runner reads \texttt{Sample\_N.sc.k}\\ and differences it term by term};
  \node[outnode, right=of f5] (ob) {\texttt{boundary\_compare.dat}\\
    \texttt{Sample\_N\_periodic.out}\\ \texttt{Sample\_N\_aperiodic.out}};
  \node[outnode, below=of f5] (o1) {\texttt{Sample\_N.out}\\
    \texttt{Sample\_N\_msg\_solid.out}\\ \texttt{Sample\_N\_compare.dat}};
  \draw[capedge] (i1) -- (f1);
  \draw[capedge] (i2) -- (f1);
  \draw[capedge] (f1) -- (f2);
  \draw[capedge] (f2) -- (p1);
  \draw[capedge] (p1) -- (oa);
  \draw[capedge] (f2) -- (f3);
  \draw[capedge] (f3) -- (f4);
  \draw[capedge] (f4) -- (f5);
  \draw[capedge] (f5) -- (o1);
  \draw[capedge2] (f4) -- (fb);
  \draw[capedge2] (fb) -- (ob);
}

\subsection{The samples and the runners}

Two TPMS samples are supplied as SwiftComp \texttt{.sc} meshes of linear tets on
a unit cell, in aluminium. Reading the \texttt{.sc} meta line directly gives
their sizes.

\begin{table}[htbp]
\centering
\small
\begin{tabular}{llrrr}
\toprule
sample & surface & nodes & tets & relative density \\
\midrule
\texttt{Sample\_1} & Schwarz-P & 116\,851 & 545\,741 & 0.3000 \\
\texttt{Sample\_2} & Schwarz-P & 54\,874 & 191\,957 & 0.0857 \\
\bottomrule
\end{tabular}
\caption{The two TPMS solid structure genes. The relative density is the
fraction of the cell filled by aluminium; it is \emph{not} the SG measure. Both
genes sit on a unit cell, so for both $\omega = 1$.}
\label{tab:solid-tpms}
\end{table}

\begin{figure}[htbp]
\centering
\includegraphics[width=0.6\linewidth]{\FIGDIR/schwarz_p_mesh.png}
\caption{The Schwarz-P triply-periodic minimal surface on its unit cell, drawn
from the real mesh of \texttt{schwarz\_p\_3Dshell.yaml} by
\texttt{examples/OpenSG\_shell/4\_get\_solid\_props\_from\_shell\_3D\_SG/}
\texttt{schwarz\_solid\_props.py}, with the SG coordinates labelled $y_1$, $y_2$
and $y_3$. This is the surface itself, meshed as a \texttt{msg-shell} SG; the
two solid genes of Table~\ref{tab:solid-tpms} discretise the same surface as a
volume of linear tets, which is why \texttt{tpms\_all\_compare.dat} can put the
two engines side by side at matched thickness.}
\label{fig:solid-schwarzp}
\end{figure}

Four runners sit in \texttt{6\_get\_solid\_props\_from\_3D\_SG}.

\begin{table}[htbp]
\centering
\footnotesize
\begin{tabular}{p{4.4cm}p{9.4cm}}
\toprule
script & what it does \\
\midrule
\texttt{sample\_1/run\_sample\_1.py}, \texttt{sample\_2/run\_sample\_2.py} &
  homogenize one sample at $E = 69$~GPa, $\nu = 0.3$ with
  \texttt{boundary="periodic"} and compare term by term against the vendor's own
  \texttt{Sample\_N.sc.k}; writes \texttt{Sample\_N.out},
  \texttt{Sample\_N\_msg\_solid.out}, \texttt{Sample\_N\_compare.dat} and
  \texttt{Sample\_N\_mesh.png} \\
\texttt{compare\_boundary\_modes.py} & runs both samples under both boundary
  treatments and writes \texttt{boundary\_compare.dat}, plus per-mode copies
  \texttt{Sample\_N\_periodic.out} and \texttt{Sample\_N\_aperiodic.out} \\
\texttt{tpms\_solid\_props.py} & the plain homogenization driver, taking the
  sample number as an optional argument \\
\texttt{sc\_to\_3dyaml.py} & converts the \texttt{.sc} tets to the OpenSG 3-D
  solid yaml and reports the tet-volume check that distinguishes a solid mesh
  from a shell surface \\
\bottomrule
\end{tabular}
\caption{The runners of example 6.}
\label{tab:solid-ex6runners}
\end{table}

\begin{lstlisting}[style=shellst]
cd examples/OpenSG-solid/6_get_solid_props_from_3D_SG/sample_1
python run_sample_1.py
\end{lstlisting}

\begin{lstlisting}[style=shellst]
cd examples/OpenSG-solid/6_get_solid_props_from_3D_SG
python compare_boundary_modes.py
\end{lstlisting}

Both scripts are kept as scripts because both do something the solver does not:
the sample runner differences the result against the vendor's \texttt{.K} and
writes \texttt{Sample\_N\_compare.dat}, and
\texttt{compare\_boundary\_modes.py} runs each sample twice and tabulates the
two treatments against each other. The homogenization on its own is one
command --- \texttt{opensg Sample\_1.yaml}, whose header already carries
\texttt{msg: solid}, \texttt{n\_model: 3} and \texttt{refined: 0}.

One detail to check before comparing numbers across scripts: the two sample
runners and \texttt{compare\_boundary\_modes.py} use $E = 69$~GPa, matched to
the SwiftComp \texttt{.K}, whereas \texttt{tpms\_solid\_props.py} and
\texttt{sc\_to\_3dyaml.py} use $E = 70$~GPa. Only the $69$~GPa runs are
comparable with the benchmark below. The \texttt{Sample\_N\_msg\_solid.out}
files that \texttt{tpms\_solid\_props.py} writes at the top level of the folder
are regenerated on each run.

\subsection{Normalization: the \texttt{.out} is already per unit cell}

For a 3-D solid SG the measure $\omega$ is the \textbf{node bounding-box
volume}, not the summed element (material) volume. Both these genes sit on a
unit cell, so $\omega = 1$ and the \texttt{.out} is written per unit cell ---
exactly the normalization a SwiftComp \texttt{.K} uses.
\texttt{Sample\_1.out} and \texttt{Sample\_1\_periodic.out} both open with
\texttt{ OpenSG msg-solid 3D elastic model, omega 1, periodic} and report
$C_{11} = 1.0190158\times10^{10}$~Pa $= 10190.1580$~MPa, against the
\texttt{1.0190158E+004}~MPa of \texttt{Sample\_1.sc.k}; Sample 2 reports
$2.3031598\times10^{9}$~Pa $= 2303.1598$~MPa against \texttt{2.3031598E+003}. The
two files are therefore \emph{directly} comparable, with no rescale in between,
and the comparison scripts' surviving \texttt{C x omega} line is a no-op at
$\omega = 1$.

This is a change from the older behaviour, in which $\omega$ was the material
volume --- $0.3$ and $0.0857$ for these two cells --- so the \texttt{.out} came
out $1/\text{relative density}$ too stiff and only the comparison table's
rescale made the numbers line up. Any note to that effect is obsolete: the
figures printed in the \texttt{.out} now are the figures in the tables below.
Every comparison table below is per unit cell, in MPa.

\subsection{The SwiftComp digit-parity benchmark}

\begin{table}[htbp]
\centering
\small
\begin{tabular}{lrrrcrrr}
\toprule
& \multicolumn{3}{c}{Sample 1} & & \multicolumn{3}{c}{Sample 2} \\
\cmidrule(lr){2-4}\cmidrule(lr){6-8}
term & msg-solid & SwiftComp & err \% & & msg-solid & SwiftComp & err \% \\
\midrule
$C_{11}$ & 10190.158 & 10190.158 & $-0.000$ & & 2303.160 & 2303.160 & $-0.000$ \\
$C_{12}$ & 5646.069 & 5646.068 & $+0.000$ & & 1612.130 & 1612.130 & $-0.000$ \\
$C_{13}$ & 5646.092 & 5646.091 & $+0.000$ & & 1612.141 & 1612.141 & $+0.000$ \\
$C_{22}$ & 10190.131 & 10190.131 & $+0.000$ & & 2303.067 & 2303.067 & $+0.000$ \\
$C_{23}$ & 5646.012 & 5646.012 & $+0.000$ & & 1612.120 & 1612.120 & $+0.000$ \\
$C_{33}$ & 10189.958 & 10189.958 & $+0.000$ & & 2303.019 & 2303.018 & $+0.000$ \\
$C_{44}$ & 4519.640 & 4519.640 & $-0.000$ & & 1106.734 & 1106.735 & $-0.000$ \\
$C_{55}$ & 4519.637 & 4519.637 & $-0.000$ & & 1106.668 & 1106.668 & $+0.000$ \\
$C_{66}$ & 4519.755 & 4519.755 & $-0.000$ & & 1106.664 & 1106.664 & $-0.000$ \\
\bottomrule
\end{tabular}
\caption{\texttt{Sample\_1\_compare.dat} and \texttt{Sample\_2\_compare.dat},
per unit cell in MPa, at $E = 69$~GPa and $\nu = 0.3$. Sample 1 is 545\,741
tets and solved in 165.9~s; Sample 2 is 191\,957 tets and solved in 61.4~s.}
\label{tab:solid-tpmsparity}
\end{table}

Digit-for-digit on both samples, to seven significant figures on every entry.
The effective law is cubic to five digits, and every symmetry-forbidden coupling
sits four to five orders below the diagonal, which is the health check on the
all-directions periodic tie. The engineering constants that the \texttt{.out}
prints for Sample 2 --- $E_1 = 9.7551133\times10^{8}$~Pa,
$G_{12} = 1.1066643\times10^{9}$~Pa, $\nu_{12} = 0.41174$ --- are per unit
cell in the same way, and only the two moduli carry that normalization;
Poisson's ratio is a ratio.

\subsection{Periodic versus aperiodic}

\texttt{compare\_boundary\_modes.py} runs the same \texttt{.sc} through both
treatments, so the percentage column measures the aperiodic boundary layer of a
single cell directly against the exact periodic answer.

\begin{table}[htbp]
\centering
\small
\begin{tabular}{lrrrcrrr}
\toprule
& \multicolumn{3}{c}{Sample 1} & & \multicolumn{3}{c}{Sample 2} \\
\cmidrule(lr){2-4}\cmidrule(lr){6-8}
term & aperiodic & periodic & \% diff & & aperiodic & periodic & \% diff \\
\midrule
$C_{11}$ & 12977.036 & 10190.158 & $+27.349$ & & 2810.329 & 2303.160 & $+22.021$ \\
$C_{22}$ & 12977.030 & 10190.131 & $+27.349$ & & 2810.363 & 2303.067 & $+22.027$ \\
$C_{33}$ & 12976.815 & 10189.958 & $+27.349$ & & 2810.251 & 2303.019 & $+22.025$ \\
$C_{12}$ & 5423.565 & 5646.069 & $-3.941$ & & 1548.627 & 1612.130 & $-3.939$ \\
$C_{13}$ & 5423.674 & 5646.092 & $-3.939$ & & 1548.954 & 1612.141 & $-3.919$ \\
$C_{23}$ & 5423.151 & 5646.012 & $-3.947$ & & 1549.049 & 1612.120 & $-3.912$ \\
$C_{44}$ & 5545.997 & 4519.640 & $+22.709$ & & 1300.352 & 1106.734 & $+17.494$ \\
$C_{55}$ & 5546.044 & 4519.637 & $+22.710$ & & 1300.151 & 1106.668 & $+17.483$ \\
$C_{66}$ & 5546.133 & 4519.755 & $+22.709$ & & 1300.048 & 1106.664 & $+17.474$ \\
\midrule
solve & 146.3~s & 150.2~s & & & 52.4~s & 51.5~s & \\
\bottomrule
\end{tabular}
\caption{\texttt{boundary\_compare.dat}, per unit cell in MPa. The aperiodic
run clamps 17\,610 boundary nodes on Sample 1 and 4456 on Sample 2.}
\label{tab:solid-boundary}
\end{table}

Read the table this way. Clamping the fluctuation on the cube faces forbids the
affine fluctuation modes, so the cell can only come out stiffer: the normal and
shear diagonals rise by $17$ to $27\,\%$, while the off-diagonal $C_{12}$ family
falls by about $4\,\%$. The bias is large because the clamped set is large ---
$17\,610$ of $116\,851$ nodes, or $15\,\%$, for Sample 1, and $4456$ of
$54\,874$, or $8\,\%$, for Sample 2 --- since the solid sheet meets the cube
faces in two-dimensional patches. On the solid route aperiodic is also
\textbf{not} faster: the periodic tie merges only face pairs, so the factorized
size barely changes, and the two timings are within a few percent of each other
in both directions. Periodic is the exact treatment for a periodic medium and is
therefore the default; aperiodic is the right model only when the SG genuinely
is not periodic.

The folder carries two further comparison tables, both regenerated by their own
scripts: \texttt{tpms\_all\_compare.dat}, which lines the solid results up
against the same Schwarz-P surface homogenized by the \texttt{msg-shell} route
at matched thickness, and \texttt{solid\_vs\_shell\_compare.dat}, which adds the
normalized stiffness $C/(E\rho)$ and the Zener ratio $2C_{44}/(C_{11}-C_{12})$
--- 1.989 for Sample 1, 3.203 for Sample 2, and 3.151 for the shell run at the
matched thickness $t = 0.0365$. Those three rows are the cleanest available
statement that the two engines describe the same structure.

\section{Beam properties from a 2-D SG}
\label{sec:solid-beam}

\subsection{What it computes}

\texttt{n\_model=1} reduces a cross-sectional SG to the Timoshenko beam
constitutive law, relating the axial force, two transverse shear forces, torque
and two bending moments to the six generalized strains
$[\epsilon_{11},\gamma_{12},\gamma_{13},\kappa_1,\kappa_2,\kappa_3]$. Internally
this is the Beam\_solid KKT engine merged into
\texttt{sg\_assembly}/\texttt{sg\_homo}: the zeroth-order warping solve $V_0$ is
run under four rigid-body Lagrange constraints, the $l$-chain first-order solve
$V_{1s}$ reuses the same factorization, and the pair is reduced to the
$6\times6$. The classical Euler-Bernoulli $4\times4$ in
$[\epsilon_{11},\kappa_1,\kappa_2,\kappa_3]$ falls out of the same $V_0$ solve
and rides along in \texttt{r["C\_eff\_EB"]} at no extra cost. The beam route
requires an SG of dimension 2 or 3.

The entry point is the same \texttt{plate\_homo\_2d} as every other route, but
\texttt{n\_model=1} branches to a different engine at the first fork, and that
engine is where the two shear rows come from.

\captree{
  \node[inputnode] (i1) {\texttt{RHC\_SW\_2UC\_45.yaml}\\
    the 2-D cross-sectional SG};
  \node[inputnode, right=of i1] (i2) {\texttt{n\_model = 1}, \texttt{name},\\
    \texttt{material\_param}, \texttt{angles}\\ in \texttt{beam\_homo\_sg.py}};
  \node[funcnode, below=of i1] (f1) {\texttt{sg\_homo.plate\_homo\_2d}\\
    $\rightarrow$ \texttt{sg\_homo.\_beam\_homo\_kkt}};
  \node[funcnode, right=of f1] (p1) {\texttt{sg\_mesh.load\_sg\_input},\\
    \texttt{sg\_mesh.plot\_sg\_mesh}};
  \node[funcnode, below=of f1] (f2)
    {\texttt{sg\_assembly.compress\_periodic\_cells\_jax},\\
     \texttt{sg\_materials.get\_heterogeneous\_C\_matrix}\\
     $\rightarrow$ \texttt{C\_asm}, \texttt{C\_stress}};
  \node[outnode, right=of f2] (oa) {\texttt{RHC\_SW\_2UC\_45\_mesh.png}};
  \node[funcnode, below=of f2] (f3)
    {\texttt{sg\_assembly.assemble\_system\_matrices},\\
     \texttt{sg\_assembly.assemble\_rigid\_body\_ops}\\
     $\rightarrow$ \texttt{Dhh} through \texttt{Dle}, \texttt{omega},\\
     \texttt{dehomo\_data}, \texttt{Psi}, \texttt{Dc}};
  \node[funcnode, below=of f3] (f4)
    {\texttt{sg\_assembly.solve\_fluctuation\_field}\\
     $\rightarrow$ \texttt{V0}, \texttt{C\_eb} under four rigid-body\\
     constraints; then \texttt{sg\_homo.prepare\_v1\_rhs} and\\
     \texttt{sg\_homo.finalize\_v1\_and\_compute\_deff}\\
     reuse the same factorization};
  \node[datanode, below=of f4] (d1) {\texttt{r["C\_eff"]} Timoshenko $6\times6$\\
    \texttt{r["C\_eff\_EB"]} classical $4\times4$\\
    \texttt{r["V0"]}, \texttt{r["V1s"]}};
  \node[funcnode, below=of d1] (f5) {\texttt{sg\_homo.write\_sc\_K}\\
    \texttt{Timoshenko} stiffness and compliance};
  \node[outnode, right=of f5] (o1) {\texttt{RHC\_SW\_2UC\_45.out}};
  \draw[capedge] (i1) -- (f1);
  \draw[capedge] (i2) -- (f1);
  \draw[capedge] (f1) -- (p1);
  \draw[capedge] (p1) -- (oa);
  \draw[capedge] (f1) -- (f2);
  \draw[capedge] (f2) -- (f3);
  \draw[capedge] (f3) -- (f4);
  \draw[capedge] (f4) -- (d1);
  \draw[capedge] (d1) -- (f5);
  \draw[capedge] (f5) -- (o1);
}

\subsection{The example SG and the input you edit}

\texttt{examples/OpenSG-solid/7\_get\_beam\_props\_from\_SG} homogenizes the same
\texttt{RHC\_SW\_2UC\_45.yaml} honeycomb cross-section as
Section~\ref{sec:solid-plate2d}: 4251 nodes, 6640 three-node triangles.
Coordinates are in mm and moduli in MPa, so the resulting stiffnesses are in N
and N$\cdot$mm$^2$.

The runner has no command-line arguments --- the SG file and the materials are
the \texttt{User Input} block at the top of the script, identical to the plate
runner except for \texttt{n\_model}.

\begin{table}[htbp]
\centering
\small
\begin{tabular}{cllll}
\toprule
id & material & $E_1, E_2, E_3$ (MPa) & $G_{12}, G_{13}, G_{23}$ (MPa) & angle \\
\midrule
1 & ply & 108000, 8000, 8000 & 4000, 4000, 3000 & $+45^\circ$ \\
2 & ply & 108000, 8000, 8000 & 4000, 4000, 3000 & $-45^\circ$ \\
3 & aluminium core & 69000, 69000, 69000 & 26540, 26540, 26540 & none \\
\bottomrule
\end{tabular}
\caption{The three materials of the honeycomb section; $\nu_{12}, \nu_{13},
\nu_{23}$ are $0.32, 0.32, 0.30$ for the plies and $0.30$ throughout for the
core.}
\label{tab:solid-beammat}
\end{table}

\begin{figure}[htbp]
\centering
\includegraphics[width=0.62\linewidth]{\FIGDIR/rhc_2dsg_mesh.png}
\caption{The mesh the beam route runs on, written by
\texttt{sg\_mesh.plot\_sg\_mesh} to \texttt{RHC\_SW\_2UC\_45\_mesh.png}. It is
the same 4251-node, 6640-triangle gene as Figure~\ref{fig:solid-rhcmesh},
because the plate and beam routes read the identical file; only
\texttt{n\_model} differs. What changes is the reading: with
\texttt{n\_model=1} the mesh is a \emph{cross-section}, its two coordinates are
$y_2$ and $y_3$, and the beam axis $y_1$ is normal to the page --- which is why
the recovery of Section~\ref{sec:solid-beamdehom} reports its extreme fibre at
$y_3 = 23.29375$~mm where the plate recovery of
Section~\ref{sec:solid-platedehom2d} reports the same point as $y_2$.}
\label{fig:solid-beammesh}
\end{figure}

\begin{lstlisting}[style=opensg,caption={beam\_homo\_sg.py, complete.}]
import numpy as np
import jax.numpy as jnp

from opensg_solid.sg_homo import plate_homo_2d

n_model = 1                    # 1: Beam; 2: Plate; 3: 3D elastic
name = "RHC_SW_2UC_45"         # reads <name>.yaml

material_param = jnp.array([
    (108e3, 8e3, 8e3, 4e3, 4e3, 3e3, 0.32, 0.32, 0.30),
    (108e3, 8e3, 8e3, 4e3, 4e3, 3e3, 0.32, 0.32, 0.30),
    (69e3, 69e3, 69e3, 26.54e3, 26.54e3, 26.54e3, 0.30, 0.30, 0.30)])
angles = jnp.array([45.0, -45.0, 0.0])

r = plate_homo_2d(name + ".yaml", material_param=material_param,
                  angles=angles, n_model=n_model)
np.set_printoptions(precision=5)
print("beam 6x6  [eps11 gam12 gam13 kappa1 kappa2 kappa3]  (Timoshenko):")
print(r["C_eff"])
print("beam 4x4  [eps11 kappa1 kappa2 kappa3]  (EB, same run):")
print(r["C_eff_EB"])
\end{lstlisting}

\subsection{Run it}

The homogenization itself is one command; the yaml in this folder carries
\texttt{n\_model: 1} and \texttt{refined: 1} in its header, so it asks for the
Timoshenko $6\times6$ without any argument:

\begin{lstlisting}[style=shellst]
cd examples/OpenSG-solid/7_get_beam_props_from_SG
opensg RHC_SW_2UC_45.yaml
\end{lstlisting}

\texttt{python beam\_homo\_sg.py} runs the same solve through the Python API
and additionally prints the classical Euler-Bernoulli $4\times4$ that rides
along in \texttt{r["C\_eff\_EB"]}, which the command-line route does not show.

\subsection{What comes out}

The run writes \texttt{RHC\_SW\_2UC\_45.out}, the timed SwiftComp-layout report
containing \texttt{The Effective Timoshenko Beam Stiffness Matrix} followed by
\texttt{The Effective Timoshenko Beam Compliance Matrix}, with banner
\texttt{ OpenSG msg-solid beam model, omega 1} and closing line
\texttt{ Time taken: 8.94 sec}; and \texttt{RHC\_SW\_2UC\_45\_mesh.png}, the SG
mesh coloured by material.

The banner reports $\omega = 1$: for a beam model on a 2-D SG the measure is
unity, so the $6\times6$ is the \textbf{total sectional stiffness}, not a
per-area density. No engineering-constants block is printed, because a beam
stiffness is not a material stiffness.

\subsection{What the numbers mean}

The shipped \texttt{RHC\_SW\_2UC\_45.out} stiffness block reads:

\begin{lstlisting}[style=shellst,basicstyle=\ttfamily\tiny]
     9.8992866E+006     6.9973096E+003     1.2577025E+001    -7.7065215E-001     2.2740495E+000    -8.6717860E-001
     6.9973096E+003     1.9074524E+006     5.9324366E+001    -3.4044009E+001    -1.6298115E+000    -6.0670378E-002
     1.2577025E+001     5.9324366E+001     9.3619028E+005    -8.4787043E+001    -9.3250231E-001     8.5948371E-001
    -7.7065184E-001    -3.4044009E+001    -8.4787043E+001     2.9774561E+008    -3.1862083E+006    -7.3057137E+003
     2.2740506E+000    -1.6298115E+000    -9.3250231E-001    -3.1862083E+006     2.1895745E+009    -9.5964309E+004
    -8.6717875E-001    -6.0670378E-002     8.5948371E-001    -7.3057137E+003    -9.5964309E+004     9.2899273E+008
\end{lstlisting}

\begin{table}[htbp]
\centering
\small
\begin{tabular}{lllp{4.6cm}}
\toprule
entry & conjugate pair & value & reads as \\
\midrule
$C_{11}$ & $\epsilon_{11}$ & $9.8992866\times10^{6}$ N & axial (extensional) stiffness \\
$C_{22}$ & $\gamma_{12}$ & $1.9074524\times10^{6}$ N & transverse shear stiffness, direction 2 \\
$C_{33}$ & $\gamma_{13}$ & $9.3619028\times10^{5}$ N & transverse shear stiffness, direction 3 \\
$C_{44}$ & $\kappa_1$ & $2.9774561\times10^{8}$ N$\cdot$mm$^2$ & torsional stiffness \\
$C_{55}$ & $\kappa_2$ & $2.1895745\times10^{9}$ N$\cdot$mm$^2$ & bending stiffness about axis 2 \\
$C_{66}$ & $\kappa_3$ & $9.2899273\times10^{8}$ N$\cdot$mm$^2$ & bending stiffness about axis 3 \\
$C_{45}$ & $\kappa_1$--$\kappa_2$ & $-3.1862083\times10^{6}$ N$\cdot$mm$^2$ & twist-bend coupling from the $\pm45^\circ$ plies \\
\bottomrule
\end{tabular}
\caption{The Timoshenko $6\times6$ of the honeycomb section, entry by entry.}
\label{tab:solid-timo}
\end{table}

Read the diagonal only as a first orientation. When couplings are present the
engineering stiffnesses are the inverses of the \textbf{compliance} diagonal,
which is why the same \texttt{.out} prints the compliance: here
$1/S_{11} = 9.89926\times10^{6}$ against $C_{11} = 9.89929\times10^{6}$
($0.0003\,\%$ apart) and $1/S_{44} = 2.97741\times10^{8}$ against
$C_{44} = 2.97746\times10^{8}$ ($0.0015\,\%$ apart), so on this section the
couplings shift the engineering values by less than $0.002\,\%$.

The twist-bend term $C_{45}$ is the physically meaningful signature of the
$\pm45^\circ$ layup. Quantify it rather than eyeballing it: at
$3.1862\times10^{6}$~N$\cdot$mm$^2$ it is a factor $93$ below $C_{44}$ and a
factor $690$ below $C_{55}$ --- two to three orders down, not three.

\subsection{The built-in cross-check: Euler-Bernoulli against Timoshenko}

\texttt{r["C\_eff\_EB"]} is the classical $4\times4$ from the same $V_0$ solve,
printed by \texttt{beam\_homo\_sg.py}. Its four entries must reproduce the
classical rows of the Timoshenko $6\times6$, since the two differ only by the
$l$-chain refinement that produces the shear terms. On this SG they agree to at
least five significant figures on every classical diagonal.

\begin{table}[htbp]
\centering
\small
\begin{tabular}{lrr}
\toprule
term & EB $4\times4$ & Timoshenko $6\times6$ \\
\midrule
extension & $9.89926177\times10^{6}$ & $9.89928657\times10^{6}$ \\
torsion & $2.97745363\times10^{8}$ & $2.97745606\times10^{8}$ \\
bending 2 & $2.18957447\times10^{9}$ & $2.18957452\times10^{9}$ \\
bending 3 & $9.28992811\times10^{8}$ & $9.28992729\times10^{8}$ \\
\bottomrule
\end{tabular}
\caption{The two classical laws from one run. The Timoshenko column is
\texttt{RHC\_SW\_2UC\_45\_beam\_Timo.out}; the EB column is the same four
numbers, recorded in \texttt{RHC\_SW\_2UC\_45\_beam\_EB.out} in example 8, a
file left over from an earlier version of that script.}
\label{tab:solid-ebvstimo}
\end{table}

The two transverse shear stiffnesses $C_{22}$ and $C_{33}$ exist \textbf{only}
in the Timoshenko $6\times6$; that is the entire reason to run
\texttt{n\_model=1} rather than read the EB block.

\section{Beam recovery}
\label{sec:solid-beamdehom}

\subsection{What it computes}

Dehomogenization takes a macro beam state and recovers the pointwise 3-D strain
and stress inside the SG, at every element Gauss point. For \texttt{n\_model=1}
this runs the generalized-Timoshenko recovery ladder
(\texttt{sg\_dehom.compute\_fluctuations\_gpu}): a successive-derivative
product-rule chain $\mathbf{F}' = \mathbf{R}\,\mathbf{F}$ seeded with
$\mathbf{F}_{1d} = \mathbf{C}_{\rm eff}^{-1}\bar{\boldsymbol{\epsilon}}$, which
builds the shear-corrected states and the zeroth- and first-order warping
displacement fields. Nothing from the homogenization is recomputed: the result
dictionary \texttt{r} carries the warping modes, the material tables and the
quadrature data, so you homogenize once and recover as often as you like.

The tree shows exactly what rides along. Everything below the dashed data node
is a contraction of quantities that already exist.

\captree{
  \node[inputnode] (i1) {\texttt{RHC\_SW\_2UC\_45.yaml}\\
    the 2-D cross-sectional SG};
  \node[funcnode, below=of i1] (f1) {\texttt{sg\_homo.plate\_homo\_2d}\\
    \texttt{n\_model=1}, the Section~\ref{sec:solid-beam} solve};
  \node[funcnode, right=of f1] (fw) {\texttt{sg\_homo.write\_sc\_K},\\
    \texttt{sg\_mesh.plot\_sg\_mesh},\\
    \texttt{np.savetxt} in the runner};
  \node[datanode, below=of f1] (d1) {\texttt{r["C\_eff"]}, \texttt{r["V0"]},
    \texttt{r["V1s"]},\\ \texttt{r["dehomo\_data"]},
    \texttt{r["C\_stress"]},\\ \texttt{r["reduced\_periodic\_cells"]},\\
    \texttt{r["elem\_rotation"]}, \texttt{r["phi\_qn"]}};
  \node[outnode, right=of d1] (o1) {\texttt{RHC\_SW\_2UC\_45.out}\\
    \texttt{RHC\_SW\_2UC\_45\_beam\_Timo.out}\\
    \texttt{RHC\_SW\_2UC\_45\_mesh.png}};
  \node[funcnode, below=of d1] (f2) {\texttt{sg\_dehom.dehom\_fields}\\
    $\rightarrow$ \texttt{sg\_dehom.compute\_fluctuations\_gpu}\\
    seeded with \texttt{C\_eff} inverse times
    \texttt{epsilon\_bar}\\ $\rightarrow$ \texttt{w\_1}, \texttt{w1s\_1},
    \texttt{w1s\_2}, \texttt{w\_2}, \texttt{st\_m\_final}};
  \node[inputnode, right=of f2] (i2) {\texttt{epsilon\_bar}\\
    \texttt{[eps11 gam12 gam13}\\ \texttt{kappa1 kappa2 kappa3]}\\
    in \texttt{beam\_dehom\_sg.py}};
  \node[funcnode, below=of f2] (f3)
    {\texttt{sg\_dehom.recover\_gauss\_point\_fields}\\
     and \texttt{sg\_dehom.\_u\_at\_gauss} $\rightarrow$\\
     \texttt{Gam}, \texttt{Sig} $(E,Q,6)$, \texttt{U} $(E,Q,3)$};
  \node[funcnode, below=of f3] (f4) {\texttt{sg\_dehom.export\_gauss}\\
    \texttt{sg\_dehom.gauss\_coords},
    \texttt{\_write\_field}};
  \node[outnode, below=of f4] (o2) {\texttt{RHC\_SW\_2UC\_45\_dehom.txt}\\
    \texttt{.vtk}, \texttt{.SM}, \texttt{.EM}, \texttt{.U}};
  \draw[capedge] (i1) -- (f1);
  \draw[capedge] (f1) -- (fw);
  \draw[capedge] (fw) -- (o1);
  \draw[capedge] (f1) -- (d1);
  \draw[capedge] (d1) -- (f2);
  \draw[capedge] (i2) -- (f2);
  \draw[capedge] (f2) -- (f3);
  \draw[capedge] (f3) -- (f4);
  \draw[capedge] (f4) -- (o2);
}

\subsection{The input you edit}

\texttt{examples/OpenSG-solid/8\_get\_beam\_dehom\_from\_SG} runs the same
honeycomb SG. Beyond the material block, the one new user input is
\texttt{epsilon\_bar}, the \textbf{macro beam strain}
$[\epsilon_{11},\gamma_{12},\gamma_{13},\kappa_1,\kappa_2,\kappa_3]$ --- a
strain six-vector, not a load vector. The shipped driver is
$\epsilon_{11} = 0.001$ combined with $\kappa_2 = 0.01$~mm$^{-1}$.

\begin{lstlisting}[style=opensg,caption={The body of beam\_dehom\_sg.py after the shared material block.}]
from opensg_solid.sg_homo import plate_homo_2d
from opensg_solid.sg_dehom import dehom_fields, export_gauss

# the MACRO beam state [eps11, gam12, gam13, kappa1, kappa2, kappa3]
epsilon_bar = jnp.array([0.001, 0.0, 0.0, 0.0, 0.01, 0.0])

r = plate_homo_2d(name + ".yaml", material_param=material_param,
                  angles=angles, n_model=n_model)
np.savetxt(name + "_beam_Timo.out", r["C_eff"], fmt="%16.8e",
           header="beam Timoshenko 6x6 [eps11 gam12 gam13 kappa1 kappa2"
                  " kappa3] (Beam_solid KKT engine) from the %dD SG %s"
                  % (r["n_sg"], name))

Gam, Sig, U = dehom_fields(r, epsilon_bar)
export_gauss(r, Gam, Sig, name + "_dehom", U_eqd=U)
for i, nm in enumerate(("xx", "yy", "zz", "yz", "xz", "xy")):
    print("  max|Sig_%s| = %.5e" % (nm, np.abs(Sig[..., i]).max()))
\end{lstlisting}

Only \textbf{classical-channel} states are valid drivers: extension, twist,
bending and their combinations. The recovery chain is kept verbatim from its
validated reference, and its seed is the compliance-scaled state
$\mathbf{C}_{\rm eff}^{-1}\bar{\boldsymbol{\epsilon}}$, which makes the two
transverse shear entries \texttt{epsilon\_bar[1]} and \texttt{epsilon\_bar[2]}
\textbf{recovery-inert} --- a shear-only driver recovers essentially zero
stress. This is a documented deviation from the refined theory stated in the
docstring of \texttt{compute\_fluctuations\_gpu}, not a switch to flip.

\subsection{Run it}

The yaml in this folder carries the same six numbers as its
\texttt{epsilon\_bar:} header key, so the whole chain --- homogenize, then
recover --- is one command with the \texttt{D} argument:

\begin{lstlisting}[style=shellst]
cd examples/OpenSG-solid/8_get_beam_dehom_from_SG
opensg RHC_SW_2UC_45.yaml D
\end{lstlisting}

\texttt{python beam\_dehom\_sg.py} is the same chain through the Python API,
and additionally writes the bare $6\times6$
\texttt{RHC\_SW\_2UC\_45\_beam\_Timo.out} through \texttt{np.savetxt}.

\subsection{What comes out}

\texttt{dehom\_fields} returns three arrays shaped $(E, Q, \cdot)$ over elements
and Gauss points: \texttt{Gam} $(E,Q,6)$ strain, \texttt{Sig} $(E,Q,6)$ stress,
and \texttt{U} $(E,Q,3)$ the \textbf{fluctuation-only} displacement (for the
beam route, the zeroth- plus first-order warping $w_1 + w_{1s}$; there is no
macro contribution in a unit-state SG recovery). \texttt{export\_gauss} writes
them out.

\begin{table}[htbp]
\centering
\footnotesize
\begin{tabular}{lp{8.2cm}}
\toprule
file & content \\
\midrule
\texttt{RHC\_SW\_2UC\_45.out} & the Timoshenko report written by
  \texttt{plate\_homo\_2d} itself, as in Section~\ref{sec:solid-beam};
  regenerated on each run \\
\texttt{RHC\_SW\_2UC\_45\_beam\_Timo.out} & the bare $6\times6$ written by the
  script through \texttt{np.savetxt}, one commented header line --- convenient
  to \texttt{np.loadtxt} \\
\texttt{RHC\_SW\_2UC\_45\_dehom.txt} & tab-separated Gauss table:
  \texttt{GP\_Index}, \texttt{Coord\_X}, \texttt{Coord\_Y}, six \texttt{Eps\_*},
  six \texttt{Sig\_*}; one header line plus 19\,920 rows (6640 triangles times
  three quadrature points) \\
\texttt{RHC\_SW\_2UC\_45\_dehom.vtk} & the same cloud as an ASCII unstructured
  grid of vertex cells, each strain and stress component a \texttt{SCALARS}
  field --- open in ParaView \\
\texttt{RHC\_SW\_2UC\_45\_dehom.SM} / \texttt{.EM} / \texttt{.U} & the OpenSG
  dehom files: stress, strain and fluctuation displacement in the
  \texttt{rm\_plate\_1D} layout, \texttt{x y z} columns then components printed
  in the order \texttt{11 22 33 12 13 23}, with the SG coordinates in the last
  \texttt{n\_sg} of the \texttt{x, y, z} slots; regenerated on each run \\
\texttt{RHC\_SW\_2UC\_45\_mesh.png} & the mesh figure \\
\bottomrule
\end{tabular}
\caption{Outputs of \texttt{beam\_dehom\_sg.py}.}
\label{tab:solid-ex8out}
\end{table}

The \texttt{.txt}/\texttt{.vtk} pair uses the SwiftComp component order
\texttt{xx yy zz yz xz xy}; the \texttt{.SM}/\texttt{.EM}/\texttt{.U} files
reorder to \texttt{11 22 33 12 13 23} on write only. Units follow the SG input
system throughout, MPa here.

\subsection{What the numbers mean}

The shipped table is the run of the driver above. Its component maxima are:

\begin{table}[htbp]
\centering
\small
\begin{tabular}{lcc}
\toprule
component & $\max\lvert\sigma\rvert$ (MPa) & $\max\lvert\varepsilon\rvert$ \\
\midrule
$xx$ & $1.529334\times10^{4}$ & $2.339375\times10^{-1}$ \\
$yy$ & $3.460157\times10^{3}$ & $2.265195\times10^{-1}$ \\
$zz$ & $5.539976\times10^{2}$ & $8.787251\times10^{-2}$ \\
$yz$ & $6.034080\times10^{2}$ & $1.344338\times10^{-1}$ \\
$xz$ & $1.176515\times10^{3}$ & $1.704428\times10^{-1}$ \\
$xy$ & $2.959409\times10^{3}$ & $1.028684\times10^{-1}$ \\
\bottomrule
\end{tabular}
\caption{Component maxima over the 19\,920 Gauss points of the shipped beam
recovery.}
\label{tab:solid-ex8max}
\end{table}

The axial field is the one to sanity-check first, because beam kinematics
predict it directly: $\epsilon_{xx} \approx \bar{\epsilon}_{11}
+ \kappa_2 y_3$ plus a warping correction. The first row of the shipped
\texttt{.txt} sits at $y_3 = 23.29375$~mm and reports
$\epsilon_{xx} = 2.339375\times10^{-1}$, against
$0.001 + 0.01 \times 23.29375 = 0.2339375$: an exact match, and the same value
is the table maximum, as it must be at the extreme fibre.

The useful content of the recovery is everything \emph{else}: the transverse and
shear components, which exist purely because the cross-section is heterogeneous
and warps, and which no beam theory can supply. The shipped driver is a
demonstration magnitude, not a service load; scale it to your own macro state
before reading the stresses as failure indices.

\begin{figure}[htbp]
\centering
\includegraphics[width=0.62\linewidth]{\FIGDIR/pv_beam_dehom_sxx.png}
\caption{ParaView render of the recovered field on the SG mesh, produced from
the \texttt{RHC\_SW\_2UC\_45\_dehom.vtk} written by this example, using
\texttt{docs/latex/render\_pv.py}. The array is \texttt{Sig\_xx}, the axial
stress $\sigma_{xx}$ in MPa (units follow the SG input, millimetres and
megapascals), on a symmetric colour scale of $\pm1.53\times10^{4}$~MPa --- the
$xx$ row of Table~\ref{tab:solid-ex8max}. Each of the 19\,920 points is one
Gauss point carried by its own vertex cell and coloured flat, so nothing is
averaged across an element boundary. The gene is the same honeycomb as
Figure~\ref{fig:solid-pvplatedehom} and the pattern is again linear in the
through-thickness coordinate, but the peak is an order of magnitude lower,
tracking this driver's extreme-fibre strain of $0.2339$ against the $2.3294$ of
the plate recovery.}
\label{fig:solid-pvbeamdehom}
\end{figure}

\section{Files at a glance}
\label{sec:solid-glance}

\begingroup\footnotesize
\begin{longtable}{cp{4.6cm}p{8.4cm}}
\caption{Every \texttt{msg-solid} capability in one table. \texttt{<name>} is
the SG file's own basename --- the \texttt{name} variable where a runner still
sets one --- and \texttt{<ff stem>} the basename of the \texttt{.ff} table.
\texttt{opensg} is the unified entry point; \texttt{opensg\_solid} and
\texttt{python -m opensg\_solid} are the same run.}
\label{tab:solid-glance} \\
\toprule
ex & you edit, then you run & it writes \\
\midrule
\endfirsthead
\toprule
ex & you edit, then you run & it writes \\
\midrule
\endhead
\bottomrule
\endfoot
1 & \texttt{layup\_db.yaml}; \texttt{python 1d\_sg.py}
  & \texttt{1dsg.yaml}, \texttt{1dsg.png},
    \texttt{layup\_db\_plate\_homo.out} \\
2 & the \texttt{USER INPUT} block; \texttt{python rm\_dehom.py}
  & \texttt{1dsg\_dehom.SM}, \texttt{.EM}, \texttt{.U}, \texttt{.vtk} \\
2 & a \texttt{.ff} table and the \texttt{dehom:} block;
    \texttt{python rm\_dehom\_ff.py ex5\_shell\_S8R\_field.ff}
  & \texttt{<ff stem>\_ff.SM}, \texttt{.EM}, \texttt{.U},
    \texttt{<ff stem>\_ff\_gauss.vtk} \\
3 & the yaml header (\texttt{n\_model}, \texttt{refined}, \texttt{analysis});
    \texttt{opensg <name>.yaml}
  & \texttt{<name>.out}, \texttt{<name>\_mesh.png} \\
4 & the above plus \texttt{epsilon\_bar:};
    \texttt{opensg <name>.yaml D}
  & \texttt{<name>.out},
    \texttt{<name>\_dehom.txt}, \texttt{.vtk}, \texttt{.SM}, \texttt{.EM},
    \texttt{.U}, \texttt{<name>\_mesh.png};
    \texttt{python plate\_dehom\_2dsg.py} adds
    \texttt{<name>\_plate\_ABD.out} \\
5 & \texttt{name} and the \texttt{paper} dictionary;
    \texttt{python solid\_homo\_2dsg.py}
  & \texttt{<name>\_Deff.out},
    \texttt{<name>\_constants\_vs\_paper.dat}, \texttt{<name>\_mesh.png} \\
6 & the \texttt{.sc} path, $E$ and $\nu$; \texttt{python run\_sample\_1.py}
    (the homogenization alone is \texttt{opensg Sample\_1.yaml})
  & \texttt{Sample\_1.out}, \texttt{Sample\_1\_msg\_solid.out},
    \texttt{Sample\_1\_compare.dat}, \texttt{Sample\_1\_mesh.png} \\
6 & nothing; \texttt{python compare\_boundary\_modes.py}
  & \texttt{boundary\_compare.dat}, \texttt{Sample\_1\_periodic.out},
    \texttt{Sample\_1\_aperiodic.out} and the Sample 2 pair \\
7 & the yaml header (\texttt{n\_model: 1}, \texttt{refined: 1});
    \texttt{opensg <name>.yaml}
  & \texttt{<name>.out}, \texttt{<name>\_mesh.png} \\
8 & the above plus \texttt{epsilon\_bar:};
    \texttt{opensg <name>.yaml D}
  & \texttt{<name>.out},
    \texttt{<name>\_dehom.txt}, \texttt{.vtk}, \texttt{.SM}, \texttt{.EM},
    \texttt{.U}, \texttt{<name>\_mesh.png}; \texttt{python beam\_dehom\_sg.py}
    adds \texttt{<name>\_beam\_Timo.out} \\
\end{longtable}
\endgroup
